\documentclass[manuscript,screen,review]{acmart}

\usepackage{amssymb}
\usepackage{amsthm}
\usepackage{enumitem}

\setcopyright{none}
\acmJournal{TOMS}

\newcommand{\llbracket}{\mathopen{[\mkern-3mu[}}
\newcommand{\rrbracket}{\mathclose{]\mkern-3mu]}}

\theoremstyle{definition}
\newtheorem{definition}{Definition}
\newtheorem{remark}{Remark}
\theoremstyle{plain}
\newtheorem{proposition}{Proposition}
\newtheorem{corollary}{Corollary}

\newcommand{\iptFilledCircle}[2]{%
  \put(#1,#2){\circle*{2}}%
}

\newcommand{\iptOpenCircle}[2]{%
  \put(#1,#2){\circle{2}}%
}

\newcommand{\iptFilledSquare}[2]{%
  \put(#1,#2){\makebox(0,0){\raisebox{-0.9mm}{\rule{1.8mm}{1.8mm}}}}%
}

\newcommand{\iptLabel}[2]{%
  \put(#1){\makebox(0,0){#2}}%
}

\newcommand{\iptRulerExample}{%
  \begingroup
  \setlength{\unitlength}{1mm}%
  \footnotesize
  \begin{picture}(86,18)
    \put(5,8){\vector(1,0){80}}
    \put(17,7){\line(0,1){2}}
    \put(33,7){\line(0,1){2}}
    \put(49,7){\line(0,1){2}}
    \put(65,7){\line(0,1){2}}
    \multiput(81,5)(0,2){3}{\line(0,1){1}}
    \iptLabel{9,4}{$0$}
    \iptLabel{17,4}{$1$}
    \iptLabel{25,4}{$2$}
    \iptLabel{33,4}{$3$}
    \iptLabel{41,4}{$4$}
    \iptLabel{49,4}{$5$}
    \iptLabel{57,4}{$6$}
    \iptLabel{65,4}{$7$}
    \iptLabel{73,4}{$8$}
    \iptLabel{17,1.5}{$b{=}1$}
    \iptLabel{81,1.5}{$e{=}9$}
    \iptOpenCircle{9}{8}
    \iptOpenCircle{25}{8}
    \iptOpenCircle{41}{8}
    \iptOpenCircle{57}{8}
    \iptOpenCircle{73}{8}
    \iptFilledCircle{17}{8}
    \iptFilledCircle{33}{8}
    \iptFilledCircle{49}{8}
    \iptFilledCircle{65}{8}
    \put(33,12){\line(1,0){16}}
    \iptLabel{41,14}{$s{=}2$}
    \iptLabel{17,11.5}{$\operatorname{glb}{=}1$}
    \iptLabel{65,11.5}{$\operatorname{lub}{=}7$}
  \end{picture}%
  \endgroup
}

\newcommand{\iptCuboidExample}{%
  \begingroup
  \setlength{\unitlength}{0.9mm}%
  \footnotesize
  \begin{picture}(114,68)
    \put(12,8){\vector(1,0){92}}
    \put(12,8){\vector(0,1){52}}
    \iptLabel{106,5}{$\dim_1$ ($R_1$)}
    \iptLabel{12,64}{$\dim_0$ ($R_0$)}
    \iptFilledCircle{12}{16}
    \iptFilledCircle{12}{28}
    \iptFilledCircle{12}{40}
    \iptFilledCircle{12}{52}
    \iptFilledCircle{24}{8}
    \iptFilledCircle{40}{8}
    \iptFilledCircle{56}{8}
    \iptFilledCircle{72}{8}
    \iptFilledCircle{88}{8}
    \iptLabel{8,16}{$1$}
    \iptLabel{8,28}{$3$}
    \iptLabel{8,40}{$5$}
    \iptLabel{8,52}{$7$}
    \iptLabel{24,4}{$-3$}
    \iptLabel{40,4}{$0$}
    \iptLabel{56,4}{$3$}
    \iptLabel{72,4}{$6$}
    \iptLabel{88,4}{$9$}
    \multiput(24,16)(0,12){4}{\iptFilledCircle{0}{0}}
    \multiput(40,16)(0,12){4}{\iptFilledCircle{0}{0}}
    \multiput(56,16)(0,12){4}{\iptFilledCircle{0}{0}}
    \multiput(72,16)(0,12){4}{\iptFilledCircle{0}{0}}
    \multiput(88,16)(0,12){4}{\iptFilledCircle{0}{0}}
    \put(21,13){\dashbox{1}(70,42){}}
    \iptLabel{26,11}{$\operatorname{glb}{=}(1,{-}3)$}
    \iptLabel{84,58}{$\operatorname{lub}{=}(7,9)$}
  \end{picture}%
  \endgroup
}

\newcommand{\iptIPTExample}{%
  \begingroup
  \setlength{\unitlength}{0.9mm}%
  \footnotesize
  \begin{picture}(112,86)
    \put(10,8){\vector(1,0){88}}
    \put(10,8){\vector(0,1){70}}
    \iptLabel{101,5}{$\dim_1$}
    \iptLabel{10,82}{$\dim_0$}
    \iptLabel{6,18}{$1$}
    \iptLabel{6,28}{$3$}
    \iptLabel{6,38}{$5$}
    \iptLabel{6,48}{$7$}
    \iptLabel{6,63}{$10$}
    \iptLabel{6,68}{$11$}
    \iptLabel{6,73}{$12$}
    \iptLabel{18,4}{$-3$}
    \iptLabel{33,4}{$0$}
    \iptLabel{43,4}{$2$}
    \iptLabel{48,4}{$3$}
    \iptLabel{53,4}{$4$}
    \iptLabel{63,4}{$6$}
    \iptLabel{73,4}{$8$}
    \iptLabel{78,4}{$9$}
    \iptLabel{83,4}{$10$}
    \iptLabel{93,4}{$12$}
    \iptFilledCircle{18}{18}
    \iptFilledCircle{18}{28}
    \iptFilledCircle{18}{38}
    \iptFilledCircle{18}{48}
    \iptFilledCircle{33}{18}
    \iptFilledCircle{33}{28}
    \iptFilledCircle{33}{38}
    \iptFilledCircle{33}{48}
    \iptFilledCircle{48}{18}
    \iptFilledCircle{48}{28}
    \iptFilledCircle{48}{38}
    \iptFilledCircle{48}{48}
    \iptFilledCircle{63}{18}
    \iptFilledCircle{63}{28}
    \iptFilledCircle{63}{38}
    \iptFilledCircle{63}{48}
    \iptFilledCircle{78}{18}
    \iptFilledCircle{78}{28}
    \iptFilledCircle{78}{38}
    \iptFilledCircle{78}{48}
    \put(15,15){\dashbox{1}(66,36){}}
    \iptLabel{85,33}{$C_0$}
    \iptFilledSquare{43}{63}
    \iptFilledSquare{43}{68}
    \iptFilledSquare{43}{73}
    \iptFilledSquare{53}{63}
    \iptFilledSquare{53}{68}
    \iptFilledSquare{53}{73}
    \iptFilledSquare{63}{63}
    \iptFilledSquare{63}{68}
    \iptFilledSquare{63}{73}
    \iptFilledSquare{73}{63}
    \iptFilledSquare{73}{68}
    \iptFilledSquare{73}{73}
    \iptFilledSquare{83}{63}
    \iptFilledSquare{83}{68}
    \iptFilledSquare{83}{73}
    \iptFilledSquare{93}{63}
    \iptFilledSquare{93}{68}
    \iptFilledSquare{93}{73}
    \put(39,59){\dashbox{1}(58,18){}}
    \iptLabel{100,68}{$C_1$}
  \end{picture}%
  \endgroup
}

\newcommand{\iptNonUniquePanelA}{%
  \setlength{\unitlength}{0.9mm}%
  \begin{picture}(62,28)
    \put(8,6){\vector(1,0){48}}
    \put(8,6){\vector(0,1){18}}
    \iptLabel{57,3}{$\dim_1$}
    \iptLabel{8,26}{$\dim_0$}
    \iptLabel{15,2}{$1$}
    \iptLabel{23,2}{$2$}
    \iptLabel{39,2}{$n{-}1$}
    \iptLabel{47,2}{$n$}
    \iptLabel{4,12}{$1$}
    \iptLabel{4,20}{$2$}
    \iptFilledCircle{15}{12}
    \iptFilledCircle{23}{12}
    \iptFilledCircle{39}{12}
    \iptFilledCircle{47}{12}
    \iptFilledCircle{23}{20}
    \iptFilledCircle{39}{20}
    \iptFilledCircle{47}{20}
    \iptLabel{31,12}{$\cdots$}
    \iptLabel{31,20}{$\cdots$}
    \put(12,9){\framebox(38,6){}}
    \put(20,17){\dashbox{1}(30,6){}}
  \end{picture}%
}

\newcommand{\iptNonUniquePanelB}{%
  \setlength{\unitlength}{0.9mm}%
  \begin{picture}(62,28)
    \put(8,6){\vector(1,0){48}}
    \put(8,6){\vector(0,1){18}}
    \iptLabel{57,3}{$\dim_1$}
    \iptLabel{8,26}{$\dim_0$}
    \iptLabel{15,2}{$1$}
    \iptLabel{23,2}{$2$}
    \iptLabel{39,2}{$n{-}1$}
    \iptLabel{47,2}{$n$}
    \iptLabel{4,12}{$1$}
    \iptLabel{4,20}{$2$}
    \iptFilledCircle{15}{12}
    \iptFilledCircle{23}{12}
    \iptFilledCircle{39}{12}
    \iptFilledCircle{47}{12}
    \iptFilledCircle{23}{20}
    \iptFilledCircle{39}{20}
    \iptFilledCircle{47}{20}
    \iptLabel{31,12}{$\cdots$}
    \iptLabel{31,20}{$\cdots$}
    \put(12,9){\dashbox{1}(6,6){}}
    \put(20,9){\framebox(30,14){}}
  \end{picture}%
}

\newcommand{\iptNonUniqueOptimal}{%
  \begingroup
  \footnotesize
  \begin{tabular}{@{}c@{\quad}c@{}}
    \iptNonUniquePanelA & \iptNonUniquePanelB
  \end{tabular}%
  \endgroup
}

\newcommand{\iptGreedyExample}{%
  \begingroup
  \setlength{\unitlength}{1mm}%
  \footnotesize
  \begin{picture}(52,44)
    \put(8,6){\vector(1,0){22}}
    \put(8,6){\vector(0,1){30}}
    \iptLabel{33,3}{$\dim_1$}
    \iptLabel{8,39}{$\dim_0$}
    \iptLabel{14,2}{$1$}
    \iptLabel{22,2}{$2$}
    \iptLabel{4,12}{$1$}
    \iptLabel{4,20}{$2$}
    \iptLabel{4,28}{$3$}
    \multiput(14,12)(0,8){3}{\iptFilledCircle{0}{0}}
    \multiput(22,12)(0,8){3}{\iptFilledCircle{0}{0}}
    \put(11,9){\dashbox{1}(14,6){}}
    \put(11,17){\dashbox{1}(14,6){}}
    \put(11,25){\dashbox{1}(14,6){}}
    \iptLabel{32,12}{$C_0$}
    \iptLabel{32,20}{$C_1$}
    \iptLabel{32,28}{$C_2$}
  \end{picture}%
  \endgroup
}

\newcommand{\iptGreedyMergeExample}{%
  \begingroup
  \setlength{\unitlength}{1mm}%
  \footnotesize
  \begin{picture}(52,44)
    \put(8,6){\vector(1,0){22}}
    \put(8,6){\vector(0,1){30}}
    \iptLabel{33,3}{$\dim_1$}
    \iptLabel{8,39}{$\dim_0$}
    \iptLabel{14,2}{$1$}
    \iptLabel{22,2}{$2$}
    \iptLabel{4,12}{$1$}
    \iptLabel{4,20}{$2$}
    \iptLabel{4,28}{$3$}
    \multiput(14,12)(0,8){3}{\iptFilledCircle{0}{0}}
    \multiput(22,12)(0,8){3}{\iptFilledCircle{0}{0}}
    \put(11,9){\dashbox{1}(14,22){}}
    \iptLabel{32,20}{$C_0$}
  \end{picture}%
  \endgroup
}

\newcommand{\iptNonOptimalPanelA}{%
  \setlength{\unitlength}{1mm}%
  \begin{picture}(46,44)
    \put(8,6){\vector(1,0){22}}
    \put(8,6){\vector(0,1){30}}
    \iptLabel{33,3}{$\dim_1$}
    \iptLabel{8,39}{$\dim_0$}
    \iptLabel{14,2}{$1$}
    \iptLabel{22,2}{$2$}
    \iptLabel{4,12}{$1$}
    \iptLabel{4,20}{$2$}
    \iptLabel{4,28}{$3$}
    \iptFilledCircle{14}{12}
    \iptFilledCircle{14}{20}
    \iptFilledCircle{22}{20}
    \iptFilledCircle{14}{28}
    \iptFilledCircle{22}{28}
    \put(11,9){\framebox(6,14){}}
    \put(19,17){\dashbox{1}(6,6){}}
    \put(11,25){\dashbox{1}(14,6){}}
    \iptLabel{32,16}{$C_0$}
    \iptLabel{32,20}{$C_1$}
    \iptLabel{32,28}{$C_2$}
  \end{picture}%
}

\newcommand{\iptNonOptimalPanelB}{%
  \setlength{\unitlength}{1mm}%
  \begin{picture}(46,44)
    \put(8,6){\vector(1,0){22}}
    \put(8,6){\vector(0,1){30}}
    \iptLabel{33,3}{$\dim_1$}
    \iptLabel{8,39}{$\dim_0$}
    \iptLabel{14,2}{$1$}
    \iptLabel{22,2}{$2$}
    \iptLabel{4,12}{$1$}
    \iptLabel{4,20}{$2$}
    \iptLabel{4,28}{$3$}
    \iptFilledCircle{14}{12}
    \iptFilledCircle{14}{20}
    \iptFilledCircle{22}{20}
    \iptFilledCircle{14}{28}
    \iptFilledCircle{22}{28}
    \put(11,9){\dashbox{1}(6,6){}}
    \put(11,17){\framebox(14,14){}}
    \iptLabel{32,12}{$C_0$}
    \iptLabel{32,24}{$C_1$}
  \end{picture}%
}

\newcommand{\iptNonOptimal}{%
  \begingroup
  \footnotesize
  \begin{tabular}{@{}c@{\qquad}c@{}}
    \iptNonOptimalPanelA & \iptNonOptimalPanelB \\
    greedy-plus-merge ($\kappa=3$) & optimal ($\kappa=2$)
  \end{tabular}%
  \endgroup
}

\newcommand{\iptPerDimLookup}{%
  \begingroup
  \footnotesize
  \setlength{\tabcolsep}{5pt}%
  \renewcommand{\arraystretch}{1.15}%
  \setlength{\unitlength}{1mm}%
  \begin{picture}(116,56)(0,8)
    \put(60,64){\makebox(0,0)[t]{%
      \begin{tabular}{@{}c@{}}
        \textit{root directory} \\
        \fbox{\begin{tabular}{lccc}
          \hline
          $\dim_0$ ruler & prefix & child & range \\
          \hline
          $\mathbf{[1,2,9)}$ & $\mathbf{0}$ & $\mathbf{A}$
            & $\mathbf{[0,20)}$ \\
          $[10,1,13)$ & $20$ & $B$ & $[20,38)$ \\
          \hline
        \end{tabular}}
      \end{tabular}}}
    \put(50,44){\vector(-2,-1){20}}
    \put(70,44){\vector(2,-1){20}}
    \put(27,31){\makebox(0,0)[t]{%
      \begin{tabular}{@{}c@{}}
        \textit{leaf $A$} \\
        \fbox{\begin{tabular}{lcc}
          \hline
          $\dim_1$ ruler & prefix & range \\
          \hline
          $\mathbf{[-3,3,12)}$ & $\mathbf{0}$ &
            $\mathbf{[0,5)}$ \\
          \hline
        \end{tabular}}
      \end{tabular}}}
    \put(93,31){\makebox(0,0)[t]{%
      \begin{tabular}{@{}c@{}}
        \textit{leaf $B$} \\
        \fbox{\begin{tabular}{lcc}
          \hline
          $\dim_1$ ruler & prefix & range \\
          \hline
          $[2,2,14)$ & $0$ & $[0,6)$ \\
          \hline
        \end{tabular}}
      \end{tabular}}}
    \put(6,16){\makebox(0,0)[lt]{%
      \begin{tabular}{@{}l@{}}
        $\operatorname{pos}_{[1,2,9)}(5)=2$, so $0 + 2 \cdot 5 = 10$ \\
        $\operatorname{pos}_{[-3,3,12)}(3)=2$, hence $\operatorname{pos}(5,3)=12$
      \end{tabular}}}
  \end{picture}%
  \endgroup
}

\newcommand{\iptCornerStealPanelA}{%
  \setlength{\unitlength}{0.9mm}%
  \begin{picture}(40,56)
    \put(8,6){\vector(1,0){18}}
    \put(8,6){\vector(0,1){42}}
    \iptLabel{29,3}{$\dim_1$}
    \iptLabel{8,51}{$\dim_0$}
    \iptLabel{14,2}{$1$}
    \iptLabel{22,2}{$2$}
    \iptLabel{4,12}{$1$}
    \iptLabel{4,20}{$2$}
    \iptLabel{4,28}{$3$}
    \iptLabel{4,36}{$4$}
    \iptLabel{4,44}{$5$}
    \iptFilledCircle{14}{12}
    \iptFilledCircle{14}{20}
    \iptFilledCircle{14}{28}
    \iptFilledCircle{14}{36}
    \iptFilledCircle{22}{36}
    \iptFilledCircle{14}{44}
    \iptFilledCircle{22}{44}
    \put(11,9){\framebox(6,22){}}
    \put(11,33){\dashbox{1}(14,14){}}
    \iptLabel{20,20}{$C_0^*$}
    \iptLabel{30,40}{$C_1^*$}
  \end{picture}%
}

\newcommand{\iptCornerStealPanelB}{%
  \setlength{\unitlength}{0.9mm}%
  \begin{picture}(40,56)
    \put(8,6){\vector(1,0){18}}
    \put(8,6){\vector(0,1){42}}
    \iptLabel{29,3}{$\dim_1$}
    \iptLabel{8,51}{$\dim_0$}
    \iptLabel{14,2}{$1$}
    \iptLabel{22,2}{$2$}
    \iptLabel{4,12}{$1$}
    \iptLabel{4,20}{$2$}
    \iptLabel{4,28}{$3$}
    \iptLabel{4,36}{$4$}
    \iptLabel{4,44}{$5$}
    \iptFilledCircle{14}{12}
    \iptFilledCircle{14}{20}
    \iptFilledCircle{14}{28}
    \iptFilledCircle{14}{36}
    \iptFilledCircle{22}{36}
    \iptFilledCircle{14}{44}
    \iptFilledCircle{22}{44}
    \put(11,9){\framebox(6,30){}}
    \put(20,34){\dashbox{1}(4,4){}}
    \put(12,42){\dashbox{1}(12,4){}}
    \put(10,32){\dashbox{1}(16,16){}}
    \iptLabel{30,24}{$C_0$}
    \iptLabel{30,36}{$C_1$}
    \iptLabel{30,44}{$C_2$}
    \iptLabel{14,36}{$\times$}
    \iptLabel{28,29}{\textit{deficit}}
    \iptLabel{28,50}{$K$}
  \end{picture}%
}

\newcommand{\iptCornerStealPanelC}{%
  \setlength{\unitlength}{0.9mm}%
  \begin{picture}(40,56)
    \put(8,6){\vector(1,0){18}}
    \put(8,6){\vector(0,1){42}}
    \iptLabel{29,3}{$\dim_1$}
    \iptLabel{8,51}{$\dim_0$}
    \iptLabel{14,2}{$1$}
    \iptLabel{22,2}{$2$}
    \iptLabel{4,12}{$1$}
    \iptLabel{4,20}{$2$}
    \iptLabel{4,28}{$3$}
    \iptLabel{4,36}{$4$}
    \iptLabel{4,44}{$5$}
    \iptFilledCircle{14}{12}
    \iptFilledCircle{14}{20}
    \iptFilledCircle{14}{28}
    \iptFilledCircle{14}{36}
    \iptFilledCircle{22}{36}
    \iptFilledCircle{14}{44}
    \iptFilledCircle{22}{44}
    \put(11,9){\framebox(6,22){}}
    \put(11,33){\dashbox{1}(14,14){}}
    \iptLabel{20,20}{$C_0^*$}
    \iptLabel{30,40}{$C_1^*$}
  \end{picture}%
}

\newcommand{\iptCornerSteal}{%
  \begingroup
  \footnotesize
  \begin{tabular}{@{}c@{\quad}c@{\quad}c@{}}
    \iptCornerStealPanelA &
      \iptCornerStealPanelB &
      \iptCornerStealPanelC \\
    (a) Optimal & (b) After steal & (c) After combine
  \end{tabular}%
  \endgroup
}

\title[Adaptive Index--Point Translator]{The Adaptive Index--Point Translator}
\subtitle{Flat Rank/Unrank for Irregular Spatial Data}

\author{Mirko Rahn}
\affiliation{%
  \institution{Fraunhofer ITWM}
  \city{Kaiserslautern}
  \country{Germany}}
\email{mirko.rahn@itwm.fraunhofer.de}

\renewcommand{\shortauthors}{Rahn}

\begin{document}

\begin{abstract}
Many high-performance applications rely on regular grids for compact
metadata and constant-cost Horner-style rank/unrank mappings between
coordinates and flat indices.
Irregular data is often regularized at the cost of
interpolation or zero-filled gaps.
We present the \emph{Adaptive Index--Point Translator} (IPT), an
adaptive, flat, and exact representation for irregular $d$-dimensional
point sets as a strictly ascending sequence of strided cuboids with
prefix sums.
A binary search identifies the relevant cuboid, and the Horner scheme
is applied inside it.
Storage for the IPT scales with the data irregularity~$\kappa$, while
index-to-point (\texttt{at}) costs $O(d + \log\kappa)$ and
point-to-index (\texttt{pos}) costs $O(d\log\kappa)$, both
collapsing to~$O(d)$ for (strided-)regular data where $\kappa=1$.
The IPT is closed
under geometric restriction, enabling processing units to store
only their sub-region's metadata.
We prove that a greedy-plus-combine construction algorithm
achieves the optimal IPT size in $O(dN)$ time on lexicographically
sorted input. Benchmarks on twenty x86 platform/compiler combinations
with regular grids, random subsets, and structured synthetic
multiple-survey layouts (all measured from sorted input) show that the
IPT reaches $169{-}180$\,Mpos/s on the tested regular grids while
staying within $1\%$ throughput of pure mixed-radix rank/unrank on
that regular case, uses $13.8\%$ of the memory of a sorted-points
table on the structured multiple-survey layouts, and accelerates
\texttt{pos} by
$9.3{-}40.7\times$ on regular grids and by ${\sim}2.33\times$ on the
structured workloads.
\end{abstract}

\ccsdesc[500]{Theory of computation~Data structures design and analysis}
\ccsdesc[300]{Software and its engineering~Data types and structures}
\ccsdesc[100]{Applied computing~Physical sciences and engineering}

\keywords{spatial indexing, rank/unrank, Horner scheme, irregular grids,
adaptive data structures, cuboid decomposition, high-performance computing}

\maketitle

\section{Introduction}
\label{section:introduction}

Many high-performance computing (HPC) applications are built around
regular grids because regularity provides simple data layouts and
constant-cost mappings between coordinates and flat indices. In
practice, regular domains are often an
approximation. Masks, embedded geometries, task-local subsets, and
partially structured domains introduce irregularity, often
gradually. Structured simulation codes often operate on sparse active
regions such as masked cells, where a regularized array is dominated
by inactive entries. In geoscience workloads, multiple acquisitions
with different extents and spacings are often forced onto a common
regular grid, increasing interpolation cost and creating large
zero-filled regions.

This paper introduces the \emph{Adaptive Index--Point Translator}
(IPT), an adaptive, flat, and exact representation for irregular
$d$-dimensional point sets that preserves $O(d)$ rank/unrank behavior
on (strided-)regular data. The domain is stored as a strictly
ascending sequence of strided cuboids. A query first locates the
relevant cuboid by binary search and then applies the Horner scheme
within that cuboid.

The IPT is part of a broader design space for extending Horner-style
rank/unrank to irregular data. At one end there is the regular-grid
case.  At the other end there is a recursive variant where each
divmod step becomes a binary search over per-dimension
directories. The IPT is in the middle: One global binary search
resolves irregularity, the inner computation remains Horner. The
advantage of the IPT is the data-independent and flat memory layout.

The IPT is useful for HPC as it preserves the (strided-)regular case
without penalty, represents data exactly and compactly, and uses a
pointer-free flat layout suitable for serialization and memory
mapping. Moreover, the IPT is closed under geometric restriction, so
each processing unit can store only its own sub-region's
metadata. These properties make it useful for local geometric
processing, geometric decomposition across processing units, and
incremental adoption in legacy regular-grid codes.

\paragraph{Effective regime.}
The empirical evaluation in this paper makes the regime in which the
IPT is competitive precise. The IPT shines on structured irregular
workloads where the cuboid count $\kappa$ stays far below the point
count $N$ and where the dominant query is
\texttt{pos}-heavy, selection-heavy, or restriction-heavy: in those
regimes it offers compact metadata, fast point-to-index lookup, and a
geometric restriction operator that closes back into the same layout.
The trade-off it pays is on \texttt{at}-heavy traversals of structured
workloads, where a sorted-points table can do less work per query;
Section~\ref{section:results-at} reports that gap honestly, and the
introduction therefore does not claim a uniform speedup. On
(strided-)regular inputs the IPT collapses to the Horner case and
costs essentially the same as a hand-written mixed-radix
implementation. On highly fragmented random subsets, where $\kappa$
approaches $N$, the structural advantage disappears and the
representation behaves like explicit point storage. The proven optimal
constructor (greedy-plus-combine) coincides exactly with the simpler
greedy-plus-merge on every one of the 61 organic workload shapes; its
practical role in our experiments is therefore correctness insurance
rather than measurable additional compression, and we frame it as such
in the rest of the paper.

Section~\ref{section:requirements} states the problem and derives the
requirements. Section~\ref{section:existing} reviews related work.
Sections~\ref{section:representation}, \ref{section:construction}, and
\ref{section:design-space} develop the representation, the construction
algorithms, and the design rationale. Sections~\ref{section:methodology}
and~\ref{section:results} present the experimental methodology and
results. Section~\ref{section:discussion} discusses limitations, and
Section~\ref{section:conclusion} concludes.

\paragraph{Contributions.}
\begin{itemize}[leftmargin=*,nosep]
  \item an adaptive flat representation for irregular $d$-dimensional point sets that
    preserves (strided-)regular-grid rank/unrank behavior,
  \item a Horner-to-IPT design space that connects the IPT to
    regular-grid indexing and clarifies the recursive-versus-flat
    trade-off,
  \item practical linear time construction algorithms (greedy-plus-merge and
    greedy-plus-combine), the latter proved optimal,
  \item complexity analyses for rank/unrank, storage, geometric
    operations, and geometric decomposability, and
  \item a benchmark implementation and empirical evaluation on twenty
    x86 platform/compiler combinations.
\end{itemize}

\section{Problem Statement and Requirements}
\label{section:requirements}

We consider a finite, non-empty set of points $D \subseteq \mathbb{Z}^d$
in fixed dimension $d$, together with a flat index space
$M = [0,N)$, $N = |D|$.
The goal is a representation of $D$ that supports efficient bidirectional translation,
\[
  \operatorname{at} : M \to D,
  \qquad
  \operatorname{pos} : D \to M,
\]
such that
\[
  \operatorname{at}(\operatorname{pos}(p)) = p
  \qquad\text{for all } p \in D,
\]
and
\[
  \operatorname{pos}(\operatorname{at}(n)) = n
  \qquad\text{for all } n \in M.
\]

For integers $b < e$, we write
\[
  [b,e) = \{x \in \mathbb{Z} \mid b \le x < e\},
\]
and for integers $b < e$ and positive stride $s$, we write
\[
  [b,s,e) = \{ b + s i \mid 0 \le i < \lceil (e-b)/s \rceil \}.
\]

The class of represented point sets must cover at least three classes:
\[
\begin{tabular}{rl}
\emph{Regular:} & $D = [b_0, e_0) \times \cdots \times [b_{d-1}, e_{d-1})$
\\
\emph{Strided-regular:} & $D = [b_0, s_0, e_0) \times \cdots \times [b_{d-1}, s_{d-1}, e_{d-1})$
\\
\emph{Irregular:} & $D \subseteq [b_0, e_0) \times \cdots \times [b_{d-1}, e_{d-1})$
\end{tabular}
\]


An important intermediate case is \emph{multiple-survey} data: the union of a few possibly overlapping strided-regular point sets.

The application requirements that motivate the IPT are:

\paragraph{Exactness.}
The representation must describe the point set exactly, without regularization by interpolation or neutral fill.

\paragraph{Adaptive space usage.}
Metadata must scale with the geometric structure present in the data,
not with the volume of a regularized bounding box. If the data
decomposes into few structured components, storage overhead must be
small, and on (strided-)regular data, metadata must be constant for fixed $d$.

\paragraph{Adaptive access time.}
Rank/unrank must be at most logarithmic in the number of structured components and reduce to constant-cost behavior on (strided-)regular inputs.

\paragraph{Flat layout.}
The representation must admit a contiguous, pointer-free in-memory layout so that it can be serialized, memory-mapped, and used directly after loading without pointer fix-up.

\paragraph{Geometric operations.}
Region-based operations such as intersection and selection must work directly on the compressed geometry, without expanding the point set.

\paragraph{Geometric decomposability.}
When data is distributed across processing units along geometric
boundaries, each unit needs to operate with only the metadata for its own sub-region.
The representation must be \emph{closed under geometric restriction}: Restricting to a sub-region must produce a valid instance.

\paragraph{Compatibility with existing regular-grid practice.}
To be a useful extension the IPT must preserve the regular-grid fast
path and therefore must not materially degrade performance when used
with (strided-)regular data.

\section{Related Work}
\label{section:existing}

This section reviews the main related approaches through the lens of the
requirements from Section~\ref{section:requirements}.

\paragraph{Explicit sparse point storage.}
Formats such as COO and the closely related CSR and CSC layouts~\cite{Saad2003,Barrett1994}
are exact and flexible because they store each point or nonzero
explicitly. They therefore satisfy exactness, but they do not compress
geometric regularity, do not provide compact rank/unrank semantics, and
do not reduce metadata cost when the point set has strong cuboid
structure.

\paragraph{Run-length and interval encodings.}
In one dimension, run-length encodings~\cite{Golomb1966} and interval
representations satisfy several of the requirements. They compress
contiguous runs, support rank/unrank, and admit flat layouts. The main
open step is the multi-dimensional extension, where one needs cuboids
in place of intervals together with both rank and unrank support.

\paragraph{Compressed bitmap indexes.}
Roaring bitmaps~\cite{ChambiLemireKaserGodin2016,LemireBoyttsovKurz2018}
encode large 32-bit set memberships by partitioning the universe into
$2^{16}$-element chunks and storing each chunk in the most compact of
three containers (sorted array, run-length, or bitmap). They are exact,
dense-friendly, and support fast intersection, but they are designed
for 1D set membership and rely on a $2^{16}$-aligned chunking that does
not generalise to multi-dimensional cuboid metadata of the form used
here. We nevertheless include the 64-bit Roaring map (CRoaring,
vendored under \path{third_party/CRoaring/}) as a measured baseline
(\texttt{Roaring} in Table~\ref{table:implementations}): each point is
linearised through the bounding grid by mixed-radix rank, the resulting
identifier is inserted into a single \texttt{Roaring64Map}, and the
range selection re-enumerates the bounding-grid cells of $\text{query}
\cap \text{grid}$ and probes membership.
This serves as the mature open-source compressed-bitmap reference for
the dense and sparse occupancy regimes in
Section~\ref{section:results-selection}.

\paragraph{Succinct rank/select dictionaries.}
Succinct data structures based on rank/select queries on bit
vectors~\cite{Jacobson1989,RamanRamanRao2007} achieve
information-theoretic optimal space for set membership while
supporting $O(1)$ rank and select. They underlie many compact spatial
indexes including $k^2$-trees for binary relations and adjacency
matrices~\cite{BrisaboaLadraNavarro2014}, which compactly encode
2D-clustered binary data. These structures are exact and compact, but
they provide rank/select on bit positions rather than directly on
multi-dimensional points, and are not pointer-free in the same flat
sense as the IPT.

\paragraph{Block-sparse and tiled sparse formats.}
Block-CSR and variable-block-row formats~\cite{ImYelickVuduc2004,Vuduc2005}
exploit clustering by grouping entries into dense sub-blocks. The ELL
and DIA formats studied in the GPU context by Bell and Garland~\cite{BellGarland2009}
follow a similar idea. These formats capture some regularity and can be
storage efficient, but they are designed around sparse linear algebra
rather than bidirectional translation between flat indices and points.

\paragraph{Higher-dimensional sparse tensor formats.}
Multi-dimensional generalisations of CSR include the Compressed Sparse
Fiber (CSF) format~\cite{SmithKarypis2015} and the hierarchical
coordinate format HiCOO~\cite{LiSunVuduc2018}. CSF stores a tensor as
a hierarchy of mode-aligned fiber arrays, while HiCOO partitions the
tensor into super-blocks and stores per-super-block local coordinates
in reduced bit width. Both compress regular structure along selected
modes, but they target sparse-tensor multiplication kernels and do not
expose a flat global rank/unrank pair over the occupied points.

\paragraph{Space-filling curves.}
Space-filling curves such as the Morton (Z-order) curve~\cite{Morton1966} and the Hilbert curve map multi-dimensional coordinates to a single linear index via bit-interleaving.
The resulting ordering preserves spatial locality well enough for range queries with B\nobreakdash-trees~\cite{LawderKing2001,TropfHerzog1981}, including run-length-compressed variants over the curve codes themselves.
However, for a sparse subset the curve still requires an explicit list
of occupied codes (or a compressed bitmap over them, which inherits
the limitations above). It therefore gives an ordering rather than a
compressed exact representation of structured sparse geometry, and its
coordinate scrambling is not compatible with Horner-style rank/unrank.

\paragraph{Spatial trees and search structures.}
Quadtrees~\cite{FinkelBentley1974}, octrees~\cite{Meagher1982},
k-d trees~\cite{Bentley1975}, and R-trees~\cite{Guttman1984} target
hierarchical spatial search. They are strong on search and pruning, but
they do not provide flat rank/unrank semantics and their pointer-based
layouts do not match the contiguous-layout requirement.

\paragraph{Hierarchical sparse grids and adaptive mesh refinement.}
Adaptive mesh refinement (AMR)~\cite{BergerOliger1984,BergerColella1989}
covers a domain with a hierarchy of axis-aligned rectangular patches,
each internally regular. This is structurally close to the IPT, but AMR
patches may overlap across refinement levels, require pointer-based
navigation, and do not provide one global flat rank/unrank map.
Hierarchical sparse volume representations such as VDB~\cite{Museth2013}
adopt a B-tree-like structure with fixed branching that uniformly
subdivides index space. Like spatial trees, they support hierarchical
sparse access rather than flat index semantics in a contiguous,
serialized layout.

\paragraph{Positioning.}
No single existing approach provides the full combination: exact
representation, adaptive compression, flat rank/unrank, contiguous
pointer-free layout, geometric decomposability, and compatibility with
regular data. The IPT representation is built on top of well-known
techniques (binary search, prefix sums, Horner rank/unrank). What is new is the specific combination that
fulfils all requirements simultaneously.

\section{IPT Representation and Core Operations}
\label{section:representation}

\subsection{Rulers}

\begin{definition}[Ruler]
  Let $b,e \in \mathbb{Z}$ with $b < e$, $s \in \mathbb{N}^+$ a
  positive stride, and $s \mid (e-b)$. Let $n = (e-b)/s \in \mathbb{N}^+$.
  A \emph{ruler} $R = [b,s,e)$ is the set
\[
R = [b,s,e) = \{ b + s i \mid i \in [0, n) \}
    = \{ b, b+s,\ldots, b+s(n-1)\}.
\]
The mappings $\operatorname{at}_R:[0, n)\to R$ and
  $\operatorname{pos}_R:R\to [0, n)$ are given by
\begin{align*}
  \operatorname{at}_R(i) &= b + s i, \text{ and}
  \\
  \operatorname{pos}_R(c) &= (c-b)/s.
\end{align*}
\end{definition}

Regular contiguous intervals $[b,e)=\{c \mid b \leq c < e\}$ are the case $s=1$, that
  is, $[b,e) = [b,1,e)$.
The cardinality of a ruler $R=[b,s,e)$ is $|R| = (e-b)/s > 0$.
Greatest lower bound and least upper bound of $R$ are $\operatorname{glb}(R) = b$ and $\operatorname{lub}(R) = e - s$.

The mapping $\operatorname{pos}_R(c)$ is well defined: $c\in R$
implies $c=b+sj$ for some $j\in[0,|R|)$, giving
  $\operatorname{pos}_R(c)= j$. It follows for every ruler $R$, index $i \in [0,|R|)$, and coordinate $c \in R$ that
$\operatorname{pos}_R(\operatorname{at}_R(i)) = i$ and $\operatorname{at}_R(\operatorname{pos}_R(c)) = c$.

\begin{figure}[ht]
\centering
\iptRulerExample
\caption{Ruler $R=[1,2,9)=\{1,3,5,7\}$. Filled circles: elements of $R$. Open circles: integers in range not in $R$. Dashed tick: exclusive end $e=9$. $\operatorname{glb}(R)=1$, $\operatorname{lub}(R)=7$, $|R|=4$.}
\Description{A one-dimensional ruler on the integers from 1 to 8. Filled points mark the elements 1, 3, 5, and 7. Open points mark excluded integers. A dashed marker indicates the exclusive end at 9, and an arrow labels the stride 2.}
\label{figure:ruler-example}
\end{figure}

\paragraph{Canonical form.}
Without the divisibility constraint, any $e'$ with $\operatorname{lub}(R) < e' \leq
\operatorname{lub}(R)+s$ would represent the same set, giving $s$
distinct representations. The condition $s \mid (e-b)$ selects the unique
representative $e = b + sn = \operatorname{lub}(R)+s$.
Several properties are then direct consequences of the definition:
(i)~$\operatorname{lub}(R) = e-s$, (ii)~$|R| = (e-b)/s$ is an exact integer operation with no rounding.
(iii)~two rulers represent the same set if and only if their triples $(b,s,e)$
are equal component-wise, enabling $O(1)$ equality.
(iv)~the adjacency condition for merging two rulers $R_1$ and $R_2$ reduces to the test $\operatorname{lub}(R_1)+s = \operatorname{glb}(R_2)$.

\subsection{Cuboids}

\begin{definition}[Cuboid]
A \emph{cuboid} is a Cartesian product of rulers,
\[
  C = R_0 \times \cdots \times R_{d-1}.
  \]
  Let $n=\prod_{j=0}^{d-1} |R_j| > 0$ and define the
  mappings $\operatorname{at}_C:[0,n)\to C$ by

\begin{align*}
  \operatorname{at}_C(i) &= (\operatorname{at}_{R_0}(q_0 \bmod
  |R_0|),\operatorname{at}_{R_1}(q_1 \bmod
  |R_1|),\dots,\operatorname{at}_{R_{d-1}}(q_{d-1} \bmod |R_{d-1}|)),
  \text{ where}
  \\ q_{d-1}&=i, \text{ and}
  \\ q_{k-1}&=\lfloor q_k / |R_k| \rfloor \quad\text{for } k=d{-}1,\ldots,1.
\end{align*}
  and $\operatorname{pos}_C:C\to [0,n)$ by
\begin{align*}
  \operatorname{pos}_C(p=(c_0,\ldots,c_{d-1})) &= i_d(p), \text{ where}
\\  i_0(p) &= 0, \text{ and}
\\  i_{k+1}(p) &= i_k(p) |R_k| + \operatorname{pos}_{R_k}(c_k) \quad\text{for } k=0,\ldots,d{-}1.
\end{align*}

\end{definition}

Regular and strided-regular data are both representable by a single cuboid.
The cardinality of $C = R_0 \times \cdots \times R_{d-1}$
is $|C| = \prod_{j=0}^{d-1} |R_j| > 0$.
Greatest lower bound and least upper bound of $C$ are
$\operatorname{glb}(C) =
(\operatorname{glb}(R_0),\ldots,\operatorname{glb}(R_{d-1}))$ and
$\operatorname{lub}(C) =
(\operatorname{lub}(R_0),\ldots,\operatorname{lub}(R_{d-1}))$.

$\operatorname{at}_C$ and $\operatorname{pos}_C$ are inverse
bijections: for every $i\in[0,|C|)$ and every $p\in C$ it holds
$
  \operatorname{pos}_C(\operatorname{at}_C(i)) = i$, and $\operatorname{at}_C(\operatorname{pos}_C(p)) = p$.
Write $n_k=|R_k|$. Every $i\in[0,|C|)$ has a unique mixed-radix expansion
\[
  i = \sum_{k=0}^{d-1} a_k \prod_{\ell=k+1}^{d-1} n_\ell,
  \qquad a_k\in[0,n_k),
\]
obtained by repeated division. The recursion for $\operatorname{at}_C$ extracts exactly these digits, $a_k=q_k\bmod n_k$, and returns the point $p=\operatorname{at}_C(i)$ with coordinates $p_k=\operatorname{at}_{R_k}(a_k)$. Applying $\operatorname{pos}_C$ to this point reinserts the same digits because $\operatorname{pos}_{R_k}(p_k)=a_k$, so the Horner recursion reconstructs the same mixed-radix value and therefore $\operatorname{pos}_C(\operatorname{at}_C(i))=i$. Conversely, for $p=(c_0,\ldots,c_{d-1})\in C$, let $a_k=\operatorname{pos}_{R_k}(c_k)$. Then $\operatorname{pos}_C(p)$ is the mixed-radix number with digits $(a_0,\ldots,a_{d-1})$, so the same uniqueness implies that the $q_k$ recursion inside $\operatorname{at}_C$ recovers those digits again. Hence each coordinate of $\operatorname{at}_C(\operatorname{pos}_C(p))$ equals $\operatorname{at}_{R_k}(a_k)=\operatorname{at}_{R_k}(\operatorname{pos}_{R_k}(c_k))=c_k$, and thus $\operatorname{at}_C(\operatorname{pos}_C(p))=p$.

\begin{figure}[ht]
\centering
\iptCuboidExample
\caption{Cuboid $C=R_0\times R_1$ where $R_0=[1,2,9)=\{1,3,5,7\}$ (the ruler of Figure~\ref{figure:ruler-example}) and $R_1=[-3,3,12)=\{-3,0,3,6,9\}$: $|R_0|=4$, $|R_1|=5$, $|C|=20$. Filled dots on the axes show the ruler elements. Dashed box marks $\operatorname{glb}(C)=(1,-3)$ to $\operatorname{lub}(C)=(7,9)$.}
\Description{A two-dimensional cuboid defined by the Cartesian product of two rulers. The x-axis ruler has elements -3, 0, 3, 6, and 9. The y-axis ruler has elements 1, 3, 5, and 7. Filled dots mark all 20 cuboid points and a dashed rectangle shows the bounding box from (1, -3) to (7, 9).}
\label{figure:cuboid-example}
\end{figure}

\subsection{Adaptive Index--Point Translator}

\begin{definition}[Adaptive Index--Point Translator]
An \emph{Adaptive Index--Point Translator} is a \emph{strictly ascending} sequence of cuboids
\[
  I = (C_0,\dots,C_{\kappa-1})
\]
where strictly ascending means
$
0 \leq i < j < \kappa \implies \operatorname{lub}(C_i) <_{\mathrm{lex}} \operatorname{glb}(C_j)
$.
Together with the sequence, an IPT stores the \emph{prefix sums of cardinalities}
\[
  r_i = \sum_{h=0}^{i-1} |C_h|, \quad r_0 = 0,
\]
so that cuboid $C_i$ occupies the global index interval
$
  [r_i, r_i + |C_i|).
$
The represented point set is
$
\llbracket I \rrbracket = \bigcup_{i=0}^{\kappa-1} C_i
$.
Let $\operatorname{ub}(x;\,v) = \min\{h : v < x_h\}$ be the \emph{upper bound} (first element strictly greater than $v$) and
$\operatorname{lb}(x;\,v) = \min\{h : v \le x_h\}$ the \emph{lower bound} (first element not less than $v$)
on a strictly increasing sequence~$x$, both standard binary-search operations.
Define the total number of points $N_I = \sum_{h=0}^{\kappa-1}|C_h|$ and the mappings $\operatorname{at}_I:[0,N_I)\to \llbracket I \rrbracket$ and $\operatorname{pos}_I:\llbracket I \rrbracket\to [0,N_I)$ by
\begin{align*}
  \operatorname{at}_I(n_I)
    &= \operatorname{at}_{C_j}(n_I-r_j),
      &&j = \operatorname{ub}(r{+}|C|;\,n_I) = \min\bigl\{h \in [0,\kappa) : n_I < r_h + |C_h|\bigr\},
  \\
  \operatorname{pos}_I(p)
    &= r_j + \operatorname{pos}_{C_j}(p),
      &&j = \operatorname{lb}(\operatorname{lub}(C);\,p) = \min\bigl\{h \in [0,\kappa) : p \le_{\mathrm{lex}} \operatorname{lub}(C_h)\bigr\}.
\end{align*}
\end{definition}

\begin{figure}[ht]
\centering
\iptIPTExample
\caption{IPT $I=(C_0,C_1)$: $C_0$ (circles, $r_0=0$, $|C_0|=20$) is the cuboid of Figure~\ref{figure:cuboid-example}. $C_1=[10,1,13)\times[2,2,14)=\{10,11,12\}\times\{2,4,6,8,10,12\}$ (squares, $r_1=20$, $|C_1|=18$). Since $\operatorname{lub}(C_0)=(7,9)<_{\mathrm{lex}}(10,2)=\operatorname{glb}(C_1)$, the sequence is strictly ascending.}
\Description{An IPT with two cuboids. The first cuboid is shown as circles on a lower grid and the second cuboid is shown as squares on a higher grid. Dashed rectangles mark the extent of each cuboid and illustrate their strict lexicographic ordering.}
\label{figure:ipt-example}
\end{figure}

Since all $|C_h|>0$, the prefix sums $r_0 < r_1 < \cdots < r_{\kappa-1}$ and the right endpoints $r_0+|C_0| < \cdots < r_{\kappa-1}+|C_{\kappa-1}|$ are strictly increasing. Strict ascent implies pairwise non-overlap and $\operatorname{lub}(C_0) <_{\mathrm{lex}} \cdots <_{\mathrm{lex}} \operatorname{lub}(C_{\kappa-1})$. Hence every $p \in \llbracket I \rrbracket$ lies in a unique cuboid $C_t$, and this index is exactly
\[
  t = \operatorname{lb}(\operatorname{lub}(C);\,p),
\]
because $\operatorname{lub}(C_h) <_{\mathrm{lex}} \operatorname{glb}(C_t) \le_{\mathrm{lex}} p$ for all $h < t$, while $p \le_{\mathrm{lex}} \operatorname{lub}(C_t)$.

Now let $n_I \in [0,N_I)$ and let
$
  j = \operatorname{ub}(r{+}|C|;\,n_I) = \min\{h \in [0,\kappa) : n_I < r_h + |C_h|\}$,
so $n_I \in [r_j, r_j + |C_j|)$. With $p = \operatorname{at}_{C_j}(n_I-r_j) \in C_j$, the characterization above gives $\operatorname{lb}(\operatorname{lub}(C);\,p)=j$, and therefore
\[
  \operatorname{pos}_I(\operatorname{at}_I(n_I)) = \operatorname{pos}_I(p)
  = r_j + \operatorname{pos}_{C_j}(p)
  = r_j + \operatorname{pos}_{C_j}(\operatorname{at}_{C_j}(n_I-r_j))
  = r_j + (n_I-r_j)
  = n_I.
\]

Conversely, let $p \in \llbracket I \rrbracket$ and let $t = \operatorname{lb}(\operatorname{lub}(C);\,p)$, so $p \in C_t$. Set $\ell = \operatorname{pos}_{C_t}(p) \in [0,|C_t|)$. Then
$
  \operatorname{pos}_I(p) = r_t + \ell \in [r_t, r_t + |C_t|)$,
and strict monotonicity of the right endpoints gives $\operatorname{ub}(r{+}|C|;\,\operatorname{pos}_I(p))=t$. Therefore
\[
  \operatorname{at}_I(\operatorname{pos}_I(p))
  = \operatorname{at}_{C_t}(\ell)
  = \operatorname{at}_{C_t}(\operatorname{pos}_{C_t}(p))
  = p.
\]
Thus for every IPT $I$, every $n_I \in [0,N_I)$, and every $p \in \llbracket I \rrbracket$,
$
  \operatorname{pos}_I(\operatorname{at}_I(n_I)) = n_I$,
  and   $\operatorname{at}_I(\operatorname{pos}_I(p)) = p
$.

If the IPT consists of one cuboid, these mappings reduce to the cuboid-local Horner scheme.

\subsection{Complexity and Storage}

\paragraph{Storage.}
The metadata cost is
\[
  \Theta(d\kappa),
\]
which becomes $\Theta(d)$ when the data is (strided-)regular ($\kappa=1$).

\paragraph{Index-to-point mapping.}
The operation \texttt{at} performs a binary search on prefix sums followed by the local cuboid mapping.
\[
  \operatorname{at} \in O(d + \log \kappa)
\]
collapsing to $O(d)$ on (strided-)regular data ($\kappa=1$).

\paragraph{Point-to-index mapping.}
The operation \texttt{pos} performs a binary search over the lexicographically ordered cuboids followed by the local cuboid mapping.
\[
  \operatorname{pos} \in O(d\log \kappa)
\]
collapsing to $O(d)$ on (strided-)regular data ($\kappa=1$).
The factor~$d$ arises because each step of the binary search compares $d$-dimensional lexicographic tuples.

\subsection{Geometric Operations and Restriction}
\label{section:operations-restriction}

\paragraph{Ruler intersection.}
Let $R=[b,s,e)$ and $R'=[b',s',e')$. Their common points are exactly the integers $c$ in the bounded interval
\[
  [\max(b,b'), \min(e,e'))
\]
that satisfy the two congruences $c \equiv b \pmod s$ and $c \equiv b' \pmod {s'}$. If these congruences are incompatible modulo $\gcd(s,s')$, the intersection is empty. Otherwise the solution set is one residue class modulo $\operatorname{lcm}(s,s')$, hence inside a bounded interval it is either empty or a canonical ruler. Computing $\gcd$, $\operatorname{lcm}$, and the first common element takes $O(1)$ arithmetic operations, so ruler intersection costs $O(1)$.

\paragraph{Cuboid intersection.}
For cuboids $C = R_0 \times \cdots \times R_{d-1}$ and $C' = R'_0 \times \cdots \times R'_{d-1}$,
\[
  C \cap C' = (R_0 \cap R'_0) \times \cdots \times (R_{d-1} \cap R'_{d-1}).
\]
Hence the intersection is empty if and only if at least one ruler intersection is empty. Otherwise, it is again a cuboid. Since there are $d$ coordinate directions and each ruler intersection costs $O(1)$, cuboid intersection costs $O(d)$.

\paragraph{IPT-cuboid intersection and selection.}
Let $Q$ be a query cuboid. A binary lower-bound search on $\operatorname{lub}(C_0),\ldots,\operatorname{lub}(C_{\kappa-1})$ finds the first cuboid whose lexicographic interval is not entirely before~$Q$. Scanning forward until $\operatorname{glb}(C_h) >_{\mathrm{lex}} \operatorname{lub}(Q)$ visits every candidate. Each candidate is tested by a cuboid intersection. If $t$ candidates are visited, the total cost is $O(d\log \kappa + dt)$.

The non-empty intersections $C_h \cap Q$ form a valid sub-IPT: each is a cuboid (by the cuboid-intersection property above), they are pairwise disjoint because the original cuboids are, and they inherit strict lexicographic ascent because every point of $C_h \cap Q$ lies inside~$C_h$. The IPT is therefore closed under geometric restriction.

These complexities match the requirements of Section~\ref{section:requirements}: storage and access scale with~$\kappa$ and collapse to $O(d)$ when $\kappa=1$.

\begin{table}[ht]
\centering
\caption{Requirements from Section~\ref{section:requirements} and the corresponding IPT properties.}
\Description{A two-column summary table mapping each stated requirement, such as exactness, adaptive space, adaptive access time, and flat layout, to the corresponding property provided by the IPT.}
\label{table:requirements-summary}
\begin{tabular}{lp{0.62\textwidth}}
\toprule
Requirement & IPT property \\
\midrule
Exactness
  & The represented set equals the union of cuboid point
    sets, with no regularization, interpolation, or fill. \\
Adaptive space
  & Metadata is $O(d\kappa)$. On (strided-)regular data $\kappa = 1$
    and the footprint is constant for fixed~$d$.
    See Section~\ref{section:results-kappa} for the measured
    structured-workload memory ratios. \\
Adaptive access time
  & \texttt{at} is $O(d+\log\kappa)$ and \texttt{pos} is
    $O(d\log\kappa)$.
    On (strided-)regular data ($\kappa = 1$) both reduce to $O(d)$. \\
Flat layout
  & A single contiguous vector of entries, with no pointers and no
    tree nodes. Directly serializable and memory-mappable. \\
Geometric operations
  & Intersection and selection operate on cuboid metadata
    without expanding the point set. \\
Geometric decomposability
  & The IPT is closed under geometric restriction. Restricting
    to a sub-region yields a valid sub-IPT formed from the
    intersected cuboids, without requiring the full global IPT. \\
Horner compatibility
  & Inside each cuboid the unchanged Horner scheme
    is reused, and the IPT adds only the cuboid-selection step. \\
\bottomrule
\end{tabular}
\end{table}

\section{Construction Algorithms}
\label{section:construction}

The IPT is only as effective as the decomposition into cuboids. This section states the decomposition problem, its complexity-theoretic baseline, and develops increasingly refined construction algorithms.

\subsection{Decomposition Problem}
\label{section:decomposition-problem}

Given a finite point set $D \subseteq \mathbb{Z}^d$, the decomposition problem is to represent $D$ as a disjoint union of cuboids,
\[
  D = C_0 \;\dot\cup\; \cdots \;\dot\cup\; C_{\kappa-1}.
\]

\paragraph{Existence and extremal size.}
Every finite point set admits an IPT decomposition, because singleton points are degenerate cuboids. Moreover,
\[
  1 \le \kappa \le N,
\]
where $\kappa=1$ if and only if the data is (strided-)regular. The upper bound $\kappa=N$ is attained when no two points belong to a common cuboid of size greater than one. For example, the set $D=\{(1,1),(2,3),(3,2)\} \subset \mathbb{Z}^2$ requires $\kappa=3$ singleton cuboids, because any pair of points differs in both dimensions with incompatible spacing, so no two of them form a $1{\times}2$, $2{\times}1$, or strided cuboid. In general, $\kappa=N$ is possible for up to $\Theta(n^{d-1})$ points in a bounding box of side length~$n$: two points that share a coordinate in $d-1$ or more dimensions always form a two-point cuboid, so every all-singletons set maps injectively into every $(d{-}1)$-dimensional projection, bounding its size by~$n^{d-1}$.

\paragraph{Non-uniqueness of optimal decompositions.}
Minimum-$\kappa$ decompositions need not be unique. As an example for $n \geq 2$, let
$
  D \;=\; \bigl([1,2){\times}[1,n{+}1)\bigr)\cup \bigl([2,3){\times}[2,n{+}1)\bigr).
$
Figure~\ref{figure:unique_decomposition_is_not_optimal} shows two different optimal decompositions of $D$.
\begin{figure}[ht]
\centering
\iptNonUniqueOptimal
\caption{Two distinct optimal decompositions of the same two-dimensional point set with minimum cuboid count $\kappa_{\mathrm{opt}}=2$.}
\Description{Two side-by-side decompositions of the same staircase-shaped point set. The left decomposition uses one long horizontal strip and one shorter strip above it. The right decomposition uses a single-point column on the left and one larger rectangle on the right. Both use two cuboids and are therefore optimal.}\label{figure:unique_decomposition_is_not_optimal}
\end{figure}
Nevertheless, because every IPT operation depends on $\kappa$ rather than on the sizes of individual cuboids, all optimal decompositions have the same storage- and query costs.

\subsection{Dynamic Programming Baseline}

Because the strictly ascending condition forces any valid decomposition to partition the lexicographically sorted input into contiguous segments, each of which forms a cuboid, the minimum-$\kappa$ decomposition can be found by dynamic programming on the sorted sequence. Let $\operatorname{opt}[i]$ denote the minimum number of cuboids needed to cover points $1, \ldots, i$. The recurrence is
\[
  \operatorname{opt}[i] = \min_{\substack{0 \leq j < i \\ \text{points } j{+}1, \ldots, i \text{ form a cuboid}}} \bigl(\operatorname{opt}[j] + 1\bigr),
\]
with $\operatorname{opt}[0] = 0$. The dynamic program considers $\Theta(N^2)$ contiguous segments. To verify that one segment forms a cuboid requires (without specific pre-processing) to scan all points in that segment, so the naive running time is $\Theta(dN^3)$ on sorted input.

While optimal, the naive $\Theta(dN^3)$ baseline (and also a potential improved $\Theta(dN^2)$ execution time) is prohibitive for large inputs. The following subsections present linear-time algorithms, culminating in greedy-plus-combine which matches the dynamic programming optimum at cost $O(dN)$.

\subsection{Greedy Construction}

Assume that the points of $D$ are given in lexicographic order. The
greedy algorithm scans the points once and maintains the current last cuboid.
\begin{enumerate}[leftmargin=*]
  \item Start the first cuboid with the first point.
  \item For each next point, test whether adding it to the last cuboid keeps that cuboid strided-regular.
  \item If yes, extend the last cuboid.
  \item Otherwise, start a new cuboid containing just that point.
\end{enumerate}

This algorithm is linear. However, even regular data can be split into multiple cuboids, because a prefix of a cuboid need not itself be a cuboid.

\paragraph{Example.}
Consider the regular grid $D = [1,4) \times [1,3)$ with $N=6$ points. In lexicographic order, the greedy algorithm processes $(1,1),(1,2),(2,1),(2,2),(3,1),(3,2)$. After $(1,1)$ and $(1,2)$, the current cuboid is $C_0 = [1,2) \times [1,3)$. The next point $(2,1)$ cannot extend $C_0$, because the set $\{(1,1),(1,2),(2,1)\}$ is not a cuboid. The algorithm therefore emits $C_0$ and starts a new cuboid. The same happens at $(3,1)$, producing three cuboids $C_0,C_1,C_2$, one per row (Figure~\ref{figure:greedy-example}), even though the input is a single regular cuboid.

\begin{figure}[ht]
\centering
\iptGreedyExample
\caption{Greedy decomposition of $D=[1,4)\times[1,3)$: three cuboids ($\kappa=3$), one per row.}
\Description{A three-by-two regular grid split by the greedy algorithm into three separate horizontal cuboids, one per row.}
\label{figure:greedy-example}
\end{figure}

\subsection{Greedy Plus Merge}

The improved algorithm augments the greedy scan with local repair.
\begin{enumerate}[leftmargin=*]
  \item Process the next point as in the basic greedy algorithm.
  \item Whenever the last cuboid changes, test whether it can be merged with the previous cuboid.
  \item If a merge is possible, merge the two cuboids and repeat the same check backward.
  \item Stop when no further backward merge is possible.
\end{enumerate}

Two cuboids can be merged if and only if they are equal in all but one dimension and their rulers in the remaining dimension extend to a single larger ruler. For non-singleton rulers the condition is the equal-stride adjacency condition, $\operatorname{lub}(R)+s = b'$ where $R$ ends at $\operatorname{lub}(R)$ with stride $s$ and the next ruler $R'$ starts at $b'$ with the same stride. Two singleton rulers always extend to a two-point ruler, and a singleton extends a non-singleton ruler exactly when it lies one stride before or after that ruler. This merge test is local and has cost $O(d)$.
The greedy mechanism maintains strict ascent: a new cuboid is only appended (or merged into the current last one) when its glb is strictly greater in lexicographic order than the lub of the current last cuboid, which is the defining condition of a valid IPT.

\paragraph{Example (continued).}
Applied to the same grid $D=[1,4)\times[1,3)$, the greedy-plus-merge algorithm initially produces $C_0=[1,2)\times[1,3)$ and $C_1=[2,3)\times[1,3)$ as before. Since $C_0$ and $C_1$ agree in dimension~1 and are adjacent in dimension~0, they merge into $C_0=[1,3)\times[1,3)$. After processing $(3,1)$ and $(3,2)$, the new $C_1=[3,4)\times[1,3)$ merges with $C_0$ in the same way, recovering the single cuboid $C_0=[1,4)\times[1,3)$ (Figure~\ref{figure:greedy-merge-example}).

\begin{figure}[ht]
\centering
\iptGreedyMergeExample
\caption{Greedy-plus-merge result for $D=[1,4)\times[1,3)$: a single cuboid ($\kappa=1$).}
\Description{The same three-by-two regular grid as in the previous figure, now shown as one merged cuboid covering all points.}
\label{figure:greedy-merge-example}
\end{figure}

\paragraph{Construction complexity.}
The greedy-plus-merge algorithm runs in time $O(dN)$ on sorted
input. Each of the $N$ greedy steps costs $O(d)$ and creates at most
one new cuboid, so at most $N$ cuboids are created in total. Each
merge reduces the cuboid count by one at cost $O(d)$, so the total number of merges is at most $N-1$. The total work is therefore $O(dN)$ for greedy steps plus $O(dN)$ for all merges.

\paragraph{Recovery of one-cuboid inputs.}
If the input point set is representable by a single cuboid, then the greedy-plus-merge algorithm reconstructs a one-cuboid decomposition. Moreover, the number of intermediate cuboids at any point during the algorithm is at most~$d$.

The proof (Appendix~\ref{proof:one-cuboid-recovery}) is by induction on~$d$. The key observation is that a $d$-dimensional cuboid decomposes in lexicographic order into constant-$\dim_0$ slices, each a $(d{-}1)$-dimensional cuboid.  By the induction hypothesis each slice collapses to a single cuboid, and consecutive slices share the same $(d{-}1)$-dimensional ruler and differ only in~$\dim_0$, so they merge.  The bound of~$d$ intermediates follows because at most one accumulated cuboid coexists with the at most $d{-}1$ intermediates of the current slice.

\paragraph{Non-optimality.}
The greedy-plus-merge algorithm does not necessarily produce a
decomposition with the minimum number of cuboids. Consider the five-point set in Figure~\ref{figure:non-optimal}
\[
  D = \{(1,1),(2,1),(2,2),(3,1),(3,2)\} \subset \mathbb{Z}^2.
\]
In lexicographic order, the greedy scan starts $C_0=[1,2){\times}[1,2)$ from $(1,1)$. At~$(2,1)$ the greedy extends~$C_0$ to $[1,3){\times}[1,2)$, because $\{(1,1),(2,1)\}$ is a valid cuboid. Then $(2,2)$ cannot extend~$C_0$, producing $C_1=[2,3){\times}[2,3)$, which does not merge with~$C_0$ (both rulers differ). Similarly, $(3,1)$ starts $C_2=[3,4){\times}[1,2)$, and $(3,2)$ extends it to $[3,4){\times}[1,3)$, which does not merge with~$C_1$. The result has $\kappa=3$ cuboids:
\[
  [1,3){\times}[1,2),\;
  [2,3){\times}[2,3),\;
  [3,4){\times}[1,3).
\]
However, a decomposition with $\kappa=2$ cuboids exists:
\[
  [1,2){\times}[1,2),\;
  [2,4){\times}[1,3).
\]

The structural cause is that the greedy criterion is purely local: at point $(2,1)$ the extension of $C_0$ from $[1,2){\times}[1,2)$ to $[1,3){\times}[1,2)$ is valid but prevents $(2,1)$ from later participating in the larger cuboid $[2,4){\times}[1,3)$. No backward merge can recover from this: neither $(C_0, C_1)$ nor $(C_1, C_2)$ satisfies the merge condition (both ruler pairs differ in each case).

\begin{figure}[ht]
\centering
\iptNonOptimal
\caption{Non-optimality of greedy-plus-merge on $D=\{(1,1),(2,1),(2,2),(3,1),(3,2)\}$. Left: the greedy-plus-merge result with $\kappa=3$. Right: an optimal decomposition with $\kappa=2$.}
\Description{Two side-by-side decompositions of a five-point set in two dimensions. The left panel shows the greedy-plus-merge result with three cuboids. The right panel shows an optimal decomposition with two cuboids.}
\label{figure:non-optimal}
\end{figure}

\paragraph{Approximation ratio.}
The corner-steal construction shows that the greedy-plus-merge can produce at least $(d+1)/2$ times as many cuboids as the optimum.

\emph{Lower bound (corner-steal).}
Let $p = (0,1,\ldots,1)$ and consider $D = \{p\} \cup [1,m{+}1)^d$ for some $m \geq 2$, with optimal decomposition $C^*_0 = \{p\}$, $C^*_1 = [1,m{+}1)^d$ ($\kappa_{\text{opt}} = 2$). The greedy scan extends $\{p\}$ with the lexicographically first corner $(1,\ldots,1)$ of~$C^*_1$ into the cuboid $[0,2){\times}[1,2)^{d-1}$. We call this a \emph{corner-steal}. The greedy-plus-merge then processes the remaining $m^d - 1$ points of $C^*_1$ and produces exactly $d$ additional cuboids, shown by induction on~$d$.

For $d = 2$, the greedy produces a partial first row
$
  [1,2){\times}[2,m{+}1)
$
and full remaining rows
$  [2,m{+}1){\times}[1,m{+}1)$.
Backward merging between them fails because their dim$_1$ rulers differ,
namely $[2,m{+}1)$ versus $[1,m{+}1)$, so the result has $2$ cuboids.

For $d \geq 3$, the first dim$_0$-slice $\{1\}{\times}[1,m{+}1)^{d-1}$
has its corner stolen and by the induction hypothesis produces $d - 1$
cuboids. The next full dim$_0$-slice
$
  \{2\}{\times}[1,m{+}1)^{d-1}
$
cannot merge with the last of those $d - 1$ cuboids because their dim$_1$
rulers differ, with $[2,m{+}1)$ in the leftover and $[1,m{+}1)$ in the
full slice. Each further full slice merges with its predecessor, so all
full slices collapse into
$[2,m{+}1){\times}[1,m{+}1)^{d-1}$
which contributes one additional cuboid. The total is $(d-1)+1 = d$.

Including the corner-steal cuboid, the greedy-plus-merge produces $d + 1$ cuboids where the optimum uses $2$, giving ratio $(d+1)/2$. By spacing $m$ independent copies in the leading dimension, the ratio persists for arbitrarily large inputs: $\kappa_{\text{gpm}} = m(d+1)$ versus $\kappa_{\text{opt}} = 2m$.

\subsection{Greedy Plus Combine}

The greedy-plus-merge algorithm can produce a suboptimal decomposition when the greedy scan absorbs a point into a narrow cuboid, preventing that point from participating in a larger cuboid formed later. The \emph{greedy-plus-combine} algorithm keeps the same greedy scan and backward merge sweep, and adds one further local repair.

\paragraph{Combine step.}
Assume the backward merge sweep has stabilised and at least three cuboids exist. Inspect the last three cuboids $C_{L-2}, C_{L-1}, C_L$. Let
\[
  K = \operatorname{bbox}(C_{L-1} \cup C_L)
\]
be the smallest cuboid containing $C_{L-1} \cup C_L$, and let
\[
  \Delta = K \setminus (C_{L-1} \cup C_L)
\]
be the missing points of that cuboid. If $\Delta \subseteq C_{L-2}$ and $R = C_{L-2} \setminus K$ is a cuboid, replace $(C_{L-2}, C_{L-1}, C_L)$ by $(R, K)$. After a successful replacement, restart the backward merge sweep and then test again for a combine. If fewer than three cuboids exist, or if either test fails, no combine is applied.

\paragraph{Why the replacement is valid.}
The replacement preserves the represented point set because $K$ already contains $C_{L-1} \cup C_L$, and its only additional points are the deficit points $\Delta$, which are taken from $C_{L-2}$. It also preserves strict ascent. The cuboid $R$ retains only points of $C_{L-2}$ that are lexicographically before $\operatorname{glb}(C_{L-1})$, while every point moved into $K$ lies at or after that boundary. $R$ is non-empty whenever the combine fires: if $C_{L-2} \subseteq K$ then $|\Delta| = |C_{L-2} \cap K| = |C_{L-2}|$, so $R = C_{L-2} \setminus K = \emptyset$, and the empty set is not a cuboid, so the combine condition~`$R$ is a cuboid' fails immediately.

\paragraph{Example (continued).}
On $D=\{(1,1),(2,1),(2,2),(3,1),(3,2)\}$, the greedy-plus-merge produces $C_0=[1,3){\times}[1,2)$, $C_1=[2,3){\times}[2,3)$, $C_2=[3,4){\times}[1,3)$. The combine step inspects all three: $K = \operatorname{bbox}(C_1 \cup C_2) = [2,4){\times}[1,3)$. The deficit is $K \setminus (C_1 \cup C_2) = \{(2,1)\}$, which is contained in $C_0$. The remainder $C_0 \setminus K = [1,2){\times}[1,2)$ is a cuboid. The algorithm replaces the three cuboids by $[1,2){\times}[1,2)$ and $[2,4){\times}[1,3)$, recovering the optimal $\kappa=2$ decomposition (Figure~\ref{figure:non-optimal}, right).

\paragraph{Construction complexity.}
Each combine attempt costs $O(d)$. Constructing $K$ and computing $|K|$ use only the coordinate bounds of the two cuboids $C_{L-1}$ and $C_L$. Since IPT cuboids are pairwise disjoint, the deficit size is
\[
  |\Delta| = |K| - |C_{L-1}| - |C_L|.
\]
Therefore $\Delta \subseteq C_{L-2}$ if and only if $C_{L-2} \cap K$ has exactly this size, which is one $O(d)$ cuboid-intersection test:
\[
  |C_{L-2} \cap K| = |K| - |C_{L-1}| - |C_L|.
\]
The remaining check that $R = C_{L-2} \setminus K$ is a cuboid is again a local test on two cuboid descriptors. A successful combine reduces the cuboid count by at least one, so at most $N-1$ combines occur over the full run. Combined with the greedy and merge costs, the greedy-plus-combine algorithm runs in $O(dN)$.

\paragraph{Optimality.}
The greedy-plus-combine always produces a decomposition with the minimum number of cuboids.
Appendix~\ref{proof:gpc-optimality} gives the full proof by strong induction on~$N$ with a secondary induction on~$d$. The key point is that the only obstruction to greedy-plus-merge is the corner-steal pattern from the previous subsection. In that configuration, the last three cuboids have exactly the shape targeted by the combine step: the missing points of $K$ lie in the previous cuboid, and removing them from that previous cuboid leaves another cuboid. The combine therefore reconstructs the first optimal cuboid~$C^*_0$. Once~$C^*_0$ has been recovered, the proof shows that no later merge or combine can modify it, so the remaining instance is smaller and the induction hypothesis applies.

\paragraph{Canonical decomposition.}
Since the minimum-$\kappa$ decomposition is not always unique (see the non-uniqueness example in~\S\ref{section:decomposition-problem}), it is natural to ask which optimum the greedy-plus-combine selects.  The inductive proof shows that the algorithm maximises $|C_0|$ first, then maximises $|C_1|$ among all optima with that $|C_0|$, and so on.  In other words, the algorithm produces the unique minimum-$\kappa$ decomposition $(C_0, \ldots, C_{\kappa-1})$ whose size sequence $(|C_0|, |C_1|, \ldots, |C_{\kappa-1}|)$ is lexicographically greatest.  We call this the \emph{canonical} optimal decomposition of~$D$.

\subsection{Comparison and Practical Choice}

Table~\ref{table:algorithm-comparison} summarises the three construction algorithms.

\begin{table}[ht]
\centering
\begin{tabular}{lccc}
\toprule
 & Greedy+Merge & Greedy+Combine & Dynamic Programming \\
\midrule
Time (sorted input) & $O(dN)$ & $O(dN)$ & $\Theta(dN^2)$, naive $\Theta(dN^3)$ \\
Optimal $\kappa$? & No & Yes & Yes \\
Approximation ratio & $\geq (d{+}1)/2$ & $1$ & $1$ \\
Working memory & ${\leq}\,d$ cuboids & ${\leq}\,d$ cuboids & $N{+}1$ sizes \\
               & (${\leq}\,3d^2$ scalars) & (${\leq}\,3d^2$ scalars) & \\
\bottomrule
\end{tabular}
\caption{Comparison of the three decomposition construction algorithms.  The greedy-plus-merge reconstructs any single cuboid (Appendix~\ref{proof:one-cuboid-recovery}) but can produce $(d{+}1)/2$ times the optimum due to corner-steals.  The greedy-plus-combine matches the dynamic programming optimum at linear cost by repairing every corner-steal via the combine step (Appendix~\ref{proof:gpc-optimality}).}
\Description{A comparison table for greedy-plus-merge, greedy-plus-combine, and dynamic programming. It lists time complexity, optimality of the cuboid count, approximation ratio, and working-memory requirements.}
\label{table:algorithm-comparison}
\end{table}

The greedy-plus-merge reconstructs any single cuboid in $O(dN)$ time with at most~$d$ intermediate cuboids (Appendix~\ref{proof:one-cuboid-recovery}).  However, the greedy scan can absorb the lexicographically first corner of a subsequent cuboid into a narrow column (the corner-steal) producing at least $(d{+}1)/2$ times the optimum, and no upper bound is currently known.

The greedy-plus-combine augments the merge with a three-cuboid combine step that detects and repairs every such corner-steal.  The formal proof (Appendix~\ref{proof:gpc-optimality}) shows by induction on~$N$ that the algorithm's first output cuboid always equals the first optimal cuboid.  If the greedy steals a corner point, the combine restores the correct cuboid, and a dimension-separation argument guarantees it is never modified again.  The result is an $O(dN)$ algorithm that always achieves the minimum cuboid count.

The dynamic programming algorithm also guarantees the minimum~$\kappa$, but its cost is $\Theta(dN^2)$ (naive: $\Theta(dN^3)$), making it impractical for large inputs.

\section{Design Space and Rationale for the Flat IPT}
\label{section:design-space}

The Horner scheme solves rank/unrank for a single regular grid.
When the point set is no longer a single cuboid, two natural strategies arise:
(i)~compose a cuboid-selection step with the unchanged Horner scheme
(the flat IPT), or
(ii)~generalise each Horner divmod step into a directory
lookup (the recursive IPT).
Both coincide for $\kappa=1$ but diverge
in structure, complexity, and practical trade-offs as
$\kappa$ grows.

\subsection{Horner Scheme}

For a regular $d$-dimensional grid with extents $n_0,\dots,n_{d-1}$ (corresponding to $|R_0|,\ldots,|R_{d-1}|$ in the cuboid notation of Section~\ref{section:representation}), the admissible coordinate tuples are
\[
  (c_0,\dots,c_{d-1}) \in [0,n_0) \times \cdots \times [0,n_{d-1}).
\]
The classical Horner rank of such a coordinate tuple is
\[
  H(c_0,\dots,c_{d-1})
  = (\cdots (((0 \cdot n_0 + c_0) n_1 + c_1)n_2 + c_2) \cdots ) n_{d-1} + c_{d-1}.
\]
Equivalently, using the same recurrence structure as the cuboid $\operatorname{pos}_C$ definition, define the partial ranks
\[
  i_0 = 0,
  \qquad
  i_{k+1} = i_k \cdot n_k + c_k
  \quad\text{for } k = 0,\ldots,d{-}1,
\]
then $H(c_0,\dots,c_{d-1}) = i_d$.

The inverse mapping, that is, the Horner unrank, reconstructs the coordinates from a flat index $i \in [0,n_0 \cdots n_{d-1})$ by repeated divmod steps. Using the same loop variable as the cuboid $\operatorname{at}_C$ definition, set $q_{d-1} = i$ and for $k=d{-}1,d{-}2,\dots,1$ let
\[
  (q_{k-1}, c_k) = \operatorname{divmod}(q_k, n_k).
\]
Finally set $c_0 = q_0$.
Then
\[
  H^{-1}(i) = (c_0,\dots,c_{d-1}).
\]

Classical regular-grid rank/unrank resolves one dimension at a time. Each divmod step selects exactly one slot. The same extends to strided-regular data, where each physical coordinate is first mapped to its ruler position before the mixed-radix mechanism applies.

\subsection{Flat IPT}

The flat IPT is not a structural generalisation of the Horner recurrence. It \emph{composes}
a cuboid-selection stage with the unchanged Horner scheme:
\begin{enumerate}[leftmargin=*]
  \item identify the cuboid that contains the requested global index or point,
  \item apply the local Horner scheme inside that cuboid.
\end{enumerate}

For \texttt{at}, a global binary search on cuboid prefix sums identifies the unique cuboid whose index interval contains the queried index. For \texttt{pos}, a global binary search on the lexicographically ordered cuboids identifies the candidate cuboid for the queried point. Once the cuboid is known, the local mapping is exactly the same rank/unrank logic that is already used in regular-grid codes.

This explains why the IPT preserves the (strided-)regular case: with the implemented one-cuboid fast path, the outer search disappears and only the cuboid-local Horner scheme remains. It also shows that the IPT extends rather than replaces existing regular-grid practice.

\subsection{Recursive IPT}

The recursive variant is included here as a design-space landmark, not
as a benchmarked competitor. We do not implement, measure, or ship it
with the artifact. Its purpose in this paper is to make the rationale
for the flat IPT precise: by spelling out the alternative
generalisation of Horner-style rank/unrank, we can compare the two
variants analytically and explain why we chose the flat IPT.

A more explicit recursive generalization is also possible. Instead of resolving irregularity with one global search, one generalizes each Horner divmod step: at dimension $j$, a binary search over admissible ruler intervals selects the next recursive step, and the local position inside that ruler replaces the corresponding divmod digit.

In the (strided-)regular case, each directory contains exactly one interval, so the binary searches are trivial. In the irregular case, the directories express ambiguity and resolve it one dimension at a time.

The attraction is that deep tuple comparisons become scalar binary searches on progressively smaller ranges. The price is more complex metadata whose footprint depends on the directory structure at each level, not just on $\kappa$.

\paragraph{Example.}
Consider the IPT of Figure~\ref{figure:ipt-example} with $d=2$ and $\kappa=2$.
The compressed recursive directory structure stores admissible ruler intervals rather than explicit coordinate values (Figure~\ref{figure:per-dim-lookup}).
At the root level there are two non-terminal directory entries,
\[
  ([1,2,9),\,0,\,A)
  \qquad\text{and}\qquad
  ([10,1,13),\,20,\,B),
\]
where each triple consists of a $\dim_0$ ruler interval, the cumulative point count before that entry, and a child reference.
The prefix~$20$ in the second entry is the size of the first entry's block, namely $|[1,2,9)|\cdot|[-3,3,12)| = 4\cdot 5 = 20$.
The two leaf directories are terminal and therefore omit the child reference:
\[
  A:([-3,3,12),\,0),
  \qquad
  B:([2,2,14),\,0).
\]
Under this compressed layout the total directory metadata is
\[
  2\cdot(3+1+1) + 2\cdot(3+1) = 18
\]
scalars: two root entries with ruler, prefix, and child reference, plus two terminal leaf entries with ruler and prefix.
By contrast, the flat IPT under its minimal cache configuration
\texttt{c0r0e0l0} stores two cuboid records with base, stride, and
bound for each of the $d=2$ rulers, plus the prefix sum
of cardinalities, for a total of
$2\cdot(2\cdot 3 + 1) = 14$~scalars; cache flags such as the
default \texttt{c0r1e1l0} add fields per record, as discussed in
Section~\ref{section:selection}.

To evaluate $\operatorname{pos}(5,3)$ with the recursive variant:
a scalar binary search in the root directory selects the first entry
$([1,2,9),0,A)$, because $5\in[1,2,9)$.
Inside that entry, the local position of~$5$ in the selected ruler is
$\operatorname{pos}_{[1,2,9)}(5)=2$.
Since the child directory~$A$ represents $|[-3,3,12)| = 5$ points per
$\dim_0$ position, the $\dim_0$ contribution is
$0 + 2\cdot 5 = 10$.
Following the child reference to~$A$, the terminal directory contributes
$\operatorname{pos}_{[-3,3,12)}(3)=2$, giving the total index
$10+2=12$.
With the flat IPT, a single lexicographic binary search on
$\operatorname{lub}(C)$ finds
$(5,3)\le_{\mathrm{lex}}(7,9)$, yielding $j=0$.
The local Horner mapping gives
$\operatorname{pos}_{R_0}(5)\cdot|R_1|+\operatorname{pos}_{R_1}(3)
 =2\cdot 5+2=12$,
and $r_0+12=12$.
Both variants yield the same index, but the recursive lookup uses two
scalar directory lookups where the flat IPT uses one two-dimensional
lexicographic comparison followed by two ruler-position computations.

\begin{figure}[ht]
\centering
\iptPerDimLookup
\caption{Compressed recursive IPT directories for $\operatorname{pos}(5,3)$
on the IPT of Figure~\ref{figure:ipt-example}.
The root stores one interval entry per admissible $\dim_0$ ruler, each with
its cumulative prefix and child reference. The terminal leaf directories store
one interval entry per admissible $\dim_1$ ruler together with the local prefix.
Bold entries trace the lookup path, yielding index~$12$.}
\Description{A schematic view of the recursive IPT lookup for the point (5,3). A root directory chooses the first-dimension ruler, a leaf directory chooses the second-dimension ruler, and the accumulated prefixes produce the final position 12.}
\label{figure:per-dim-lookup}
\end{figure}

\begin{remark}[Membership testing]
Any robust implementation of \texttt{pos} must report when the input
point is not contained in~$D$.
In the recursive variant, non-membership is detected for
free: if at some level the queried coordinate is contained in none of
the admissible ruler intervals for the current prefix, the lookup
terminates immediately.
In the flat IPT, the global binary search always returns a candidate
cuboid.
Non-membership must therefore be detected explicitly by verifying, for
each ruler of the candidate cuboid, that the corresponding coordinate
is a member.
This adds an $O(d)$ post-check that the recursive variant avoids.
\end{remark}

\subsection{Comparison}
Recall that the flat IPT stores $\Theta(d\kappa)$ scalars and provides
$\texttt{at}\in O(d+\log\kappa)$ and
$\texttt{pos}\in O(d\log\kappa)$.
The recursive variant replaces the single global search
with $d$~scalar binary searches, one per dimension level.
At each level the directory stores the admissible ruler intervals
for that dimension given the resolved prefix.
With $\kappa$~cuboids, each directory contains at most $\kappa$~intervals.
Let $n_k$ denote the directory size at level~$k$ along a lookup path.
Both \texttt{at} and \texttt{pos} then cost
$O\!\bigl(\sum_{k=0}^{d-1}\log n_k\bigr)$
scalar comparisons.
The trade-off depends on the structure of the data.

\emph{(Strided-)regular ($\kappa=1$).}
Both variants store $\Theta(d)$ and evaluate \texttt{at} and \texttt{pos}
in~$O(d)$.
Every recursive directory contains a single ruler interval,
so the binary searches are trivial.

\emph{Highly regular ($\kappa\ll N$).}
The cuboids are large, and ruler intervals compress long arithmetic
progressions into constant-size descriptors.
Because each directory level stores at most $\kappa$~ruler intervals,
the recursive metadata is compact (bounded polynomially
in~$\kappa$ and~$d$, independently of~$N$), comparable
to the flat IPT's $\Theta(d\kappa)$.
For \texttt{pos}, the flat IPT performs $O(\log\kappa)$ lexicographic
comparisons of worst-case cost $O(d)$ each, giving $O(d\log\kappa)$.
The recursive variant also costs $O(d\log\kappa)$ in the worst case
(each of the $d$~directories has up to $\kappa$~entries),
but uses $O(1)$~scalar comparisons instead of $O(d)$~tuple comparisons.
When the cuboid projections separate early, then the remaining $d{-}1$ directories
have size one and the total drops to $O(d+\log\kappa)$.
Conversely, the flat IPT's lexicographic comparisons can also
terminate early when the leading dimension already distinguishes
the cuboids.
For \texttt{at}, the flat IPT always identifies the cuboid with
a single $O(\log\kappa)$ scalar search on prefix sums,
then extracts coordinates with $O(d)$~local divmod steps,
giving $O(d+\log\kappa)$ independently of the dimensional structure.
The recursive variant identifies the cuboid incrementally
across levels:
if an early dimension already resolves the cuboid identity,
the remaining directories have size one and the total is
$O(d+\log\kappa)$.
If, however, cuboids are not distinguishable until the last
dimension, the identification cost is spread across all levels,
giving up to $O(d\log\kappa)$ in the worst case.

\emph{Highly irregular ($\kappa\approx N$).}
Each cuboid degenerates to one or a few points, so rulers are short
and the compression advantage vanishes.
Both variants store $\Theta(dN)$.
The recursive top-level directory may list up to~$N$ entries, but deeper directories narrow quickly
(each point resides in exactly one cuboid),
so $\sum\log n_k$ remains $O(d\log N)$ in the worst case, which is the
same as the flat IPT's \texttt{pos}.
For \texttt{at}, the flat IPT retains a guaranteed $O(d+\log N)$,
while the recursive variant ranges from $O(d+\log N)$
(early separation) to $O(d\log N)$ (late separation).
Neither variant has a clear advantage for \texttt{pos}.
For \texttt{at}, the flat IPT's guaranteed bound is tighter.

\paragraph{Intersection and selection.}
Ruler intersection ($O(1)$) and cuboid intersection ($O(d)$) are
representation-independent: they operate on the mathematical objects,
not on how the IPT stores them, so both variants use the same algorithm
at the same cost.
The IPT-cuboid intersection and selection operation
is where the two representations genuinely differ.
The flat IPT performs a single binary search on $\operatorname{lub}(C)$
to find the first candidate cuboid, then scans forward testing each of
$t$~candidates via an~$O(d)$ cuboid intersection, for a total of
$O(d\log\kappa+dt)$.
The recursive variant naturally supports a \emph{recursive pruning
traversal}: at dimension level~$k$, binary-search the directory for
entries compatible with the query ruler in that dimension, then recurse
into each matching subdirectory.
Let $m_k$ denote the number of subdirectory nodes visited at level~$k$
($m_0=1$). At each node a binary search on a directory of size
at most~$\kappa$ costs $O(\log\kappa)$, and the matches become nodes
at the next level.
The total cost is
$O\!\bigl(\log\kappa\cdot\sum_{k=0}^{d-1}m_k\bigr)$.
When an early dimension already prunes most branches,
$m_k=1$ for all subsequent levels and the bound collapses
to~$O(d\log\kappa)$.
When no dimension prunes effectively, $m_k$ can reach~$\kappa$
at every level, giving up to $O(d\kappa\log\kappa)$,
which is a $\log\kappa$ factor above the flat IPT's
$O(d\log\kappa+dt)$ (with $t\le\kappa$).
The flat IPT's guarantee is tighter and more predictable.
The recursive variant's cost depends on how the query
decomposes across dimensions.

\medskip
\noindent
The summary is therefore: in the (strided-)regular case both variants
collapse to $\Theta(d)$ storage and $O(d)$ queries; in the highly
regular regime the asymptotic costs match, but the recursive variant
trades $O(d)$ tuple comparisons for $O(d)$ scalar comparisons at the
price of pointer-chasing access and a query bound that depends on how
the cuboids project across dimensions; in the highly irregular regime
the flat IPT keeps a guaranteed $O(d+\log N)$ \texttt{at} cost while
the recursive variant ranges between $O(d+\log N)$ and $O(d\log N)$
depending on early vs.\ late separation. For selection, the flat IPT
pays $O(d\log\kappa+dt)$, while the recursive traversal pays
$O(\log\kappa\cdot\sum m_k)$ on the visited directory nodes, which
gains only when $\sum m_k$ is much smaller than $dt$.

\subsection{Design Choice}

The flat IPT studied in this paper is the systems-oriented member of
the design space. The recursive variant can match or occasionally improve on
the flat IPT when the data separates early, but its bounds depend on
how the cuboids project across dimensions. By contrast, the flat IPT
provides data-independent worst-case guarantees and a contiguous array
of fixed-size cuboid records that is easy to serialize and cache.

Within the cases considered here, the clearest niche for the recursive
variant is selection queries in which many cuboids share the same
leading-dimension projections and separate only in later dimensions.
In that regime, the flat IPT pays $O(d)$ for each of the $t$ candidate
cuboids, giving $O(d\log\kappa+dt)$ overall, whereas the recursive
traversal keeps the shared prefix collapsed and pays
$O\!\bigl(\log\kappa\cdot\sum m_k\bigr)$ on the visited directory
nodes instead. The potential gain is therefore in the candidate
processing term, and it appears only when $\sum m_k$ is much smaller
than $dt$. Even there, the benefit is narrow because it is paired with
a weaker \texttt{at} guarantee and pointer-chasing access.

That the simpler two-stage composition (binary search plus Horner
scheme) yields tighter and more predictable bounds than the
structurally deeper recursive generalisation is the key insight from the design-space analysis.

\section{Experimental Methodology}
\label{section:methodology}

\subsection{Implementation}
\label{section:implementation}

All experiments use a self-contained C++23 benchmark
implementation~\cite{RahnIPT2026} with no external library dependencies.

\paragraph{Data layout.}
Coordinates are \texttt{std::int64\_t}, and flat indices are
\texttt{std::uint64\_t}.
A $d$-dimensional point is a \texttt{std::array} of $d$ coordinates.
A ruler $[b,s,e)$ stores three coordinates.
A cuboid is a \texttt{std::array} of $d$ rulers.
An \emph{entry} pairs a cuboid with the cumulative size of all preceding
entries.
The IPT is a single \texttt{std::vector} of entries in strict
lexicographic order, and all orderings use compiler-generated
\texttt{operator<=>}.

\paragraph{Construction.}
Both the greedy-plus-merge and greedy-plus-combine algorithms
(Section~\ref{section:construction}) are implemented.
Points are consumed one at a time in lexicographic order. The entry
vector is maintained incrementally with recursive merging and, in the
combine variant, a subsequent three-entry combine attempt.

\paragraph{Query operations.}
\texttt{at} performs \texttt{std::upper\_bound} on the entry vector by
\texttt{end()} to locate the containing entry, then applies the cuboid's
mixed-radix decoding.
\texttt{pos} performs \texttt{std::lower\_bound} by \texttt{lub()} to find
the first candidate entry, then computes the Horner-style index and adds
the entry's stored cumulative prefix.
Intersection is exercised inside the dedicated selection and
restriction benchmarks of
Section~\ref{section:results-selection}. Those measurements capture not
only the per-fragment translation work also used by \texttt{at} and
\texttt{pos}, but also the candidate-enumeration term of the selection
loop and the from-range rebuild used for restriction.

\paragraph{Measured implementations.}
The benchmark covers eleven exact implementations
(Table~\ref{table:implementations}). The two IPT builders are
greedy-plus-merge and greedy-plus-combine. \texttt{Regular-IPT} builds
the single-cuboid IPT for fully regular data. \texttt{Grid} applies
pure mixed-radix rank/unrank to a full regular grid. \texttt{SortedPoints}
stores the explicit lexicographically sorted point array. \texttt{DenseBitset}
and \texttt{BlockBitmap} represent the bounding grid as dense occupancy
bitmaps. \texttt{Roaring} wraps the vendored CRoaring \texttt{Roaring64Map}
around the same bounding-grid linearisation, providing a mature
open-source compressed-bitmap reference. \texttt{LexRun} stores maximal
lexicographic runs in the fastest dimension. \texttt{RowCSR-k1} and
\texttt{RowCSR-k2} are compressed-sparse-row layouts with one and two row
dimensions, respectively.
Together these baselines span explicit-point, dense-occupancy,
compressed-bitmap, run-length, row-compressed, and regular-grid designs.
They do not exhaust every practical competitor, especially tree-based or
space-filling-curve variants, but they cover a broad set of exact flat
and sparse layouts that isolate the IPT's trade-offs against nearby
representations.

\begin{table}[ht]
\centering
\caption{Measured implementations.}
\Description{A table listing the ten benchmarked implementations, their representation, where they apply, and their role in the study.}
\label{table:implementations}
\small
\begin{tabular}{p{2.8cm}p{3.2cm}p{2.3cm}p{4.0cm}}
\toprule
Implementation & Representation & Scope & Role \\
\midrule
Greedy-plus-merge
  & Full IPT via incremental merge
  & all workloads
  & practical IPT default in the detailed plots \\
Greedy-plus-combine
  & Full IPT with additional combine repair
  & all workloads
  & optimal-construction IPT variant \\
Regular-IPT
  & single-cuboid IPT
  & regular grids
  & ideal regular-case IPT baseline \\
SortedPoints
  & sorted explicit point array
  & all workloads
  & general exact baseline \\
DenseBitset
  & dense occupancy bitset on the bounding grid
  & smaller structured scenarios
  & dense-grid occupancy baseline \\
BlockBitmap
  & tiled 64-bit occupancy bitmap
  & smaller structured scenarios
  & cache-local occupancy baseline \\
Roaring
  & 64-bit Roaring bitmap on the linearised bounding grid (CRoaring)
  & smaller structured scenarios
  & mature compressed-bitmap baseline \\
LexRun
  & run-length encoding in the fastest dimension
  & all core workloads
  & one-axis compression baseline \\
RowCSR-k1
  & CSR with one row dimension
  & all core workloads
  & row-oriented sparse baseline \\
RowCSR-k2
  & CSR with two row dimensions
  & all core workloads
  & deeper row-oriented sparse baseline \\
Grid
  & pure Horner rank/unrank on the full grid
  & regular grids
  & ideal dense-grid baseline \\
\bottomrule
\end{tabular}
\end{table}

\paragraph{Benchmark vs.\ industrial implementation.}
The benchmarked implementation targets C++23, uses \texttt{std::vector}
with doubling growth, throws on out-of-range queries (no throw is
hit on any measured hot path), and exposes the selection pipeline
needed for the Section~\ref{section:results-selection} measurements:
\texttt{IPT} provides \texttt{select}, and \texttt{restrict} is
evaluated by feeding that view into the \texttt{std::from\_range}
constructor. The cuboid-search and cuboid-local translation kernels
that dominate the measured hot paths are unchanged from the deployed
C++14 implementation, which differs only in surrounding storage and
error-reporting strategy (single-allocation/mmap-friendly buffers,
\texttt{expected}-style return type).

The benchmark driver evaluates the cache-dependent IPT implementations
(greedy-plus-merge, greedy-plus-combine, and Regular-IPT) under all
16 cache configurations.
The cache-independent baselines (SortedPoints, LexRun, RowCSR-k1,
RowCSR-k2, and Grid) are executed once per platform, because their data
layout does not depend on IPT cache flags.
The three occupancy-bitmap baselines (DenseBitset, BlockBitmap, and
Roaring) are restricted to the seven smaller structured scenarios because
their selection cost scales with the bounding-grid volume rather than
with the number of occupied points.

\subsection{Benchmark Setup}
\label{section:setup}

The benchmark is compiled with
\texttt{-std=c++23 -O3 -DNDEBUG} using the operating-system-provided
compiler on each platform.
Table~\ref{table:platforms} lists the 20 platform/compiler combinations,
grouped into ten hardware/OS rows.

\begin{table}[ht]
\centering
\caption{Benchmark platforms.}
\Description{A table listing ten hardware and operating-system rows that together yield twenty benchmark platform and compiler combinations. Each row includes the processor model, operating system, kernel, and the available GCC and Clang versions.}
\label{table:platforms}
\begin{tabular}{lllll}
\toprule
Processor & OS & Kernel & GCC & Clang \\
\midrule
AMD EPYC 9454
  & Rocky 10.1 & 6.12
  & 14.3.1 & 20.1.8 \\
AMD EPYC 9454
  & Rocky 9.7 & 5.14
  & 14.2.1 & 20.1.8 \\
AMD EPYC 9454
  & Rocky 8.10 & 4.18
  & 14.2.1 & 20.1.8 \\
AMD EPYC 7543
  & Rocky 10.1 & 6.12
  & 14.3.1 & 20.1.8 \\
AMD EPYC 7543
  & Rocky 9.7 & 5.14
  & 14.2.1 & 20.1.8 \\
AMD EPYC 7543
  & Rocky 8.10 & 4.18
  & 14.2.1 & 20.1.8 \\
Intel Xeon Gold 6442Y
  & Rocky 9.7 & 5.14
  & 14.2.1 & 20.1.8 \\
Intel Xeon Gold 6240R
  & Rocky 9.7 & 5.14
  & 14.2.1 & 20.1.8 \\
Intel Xeon Gold 6348
  & RHEL 9.7 & 5.14
  & 15.2.0 & 20.1.8 \\
Intel Core i7-1370P
  & Rocky 9.7 & 7.0
  & 14.2.1 & 20.1.8 \\
\bottomrule
\end{tabular}
\end{table}

The detailed per-scenario plots use the AMD EPYC~9454 on Rocky~10.1 with
GCC~14.3.1 as the reference platform.
The cache-selection plot instead aggregates all 20 platform/compiler
combinations, because that selection argument is meant to be
cross-platform.

\subsection{Workloads}
\label{section:workloads}

Two workload families are used, totaling 61 workload shapes.
\emph{Grid scenarios:}
four 3D grids ($10{\times}10{\times}10$, $100{\times}10{\times}10$,
$100{\times}100{\times}10$, $100{\times}100{\times}100$) are each tested
in 13 occupancy patterns: the fully regular case and 12 random-subset
densities (1, 2, 5, 10, 25, 50, 75, 90, 95, 98, 99, 100\,\% of the grid
points).
This yields 52 regular/random workload shapes.
Each random case is repeated over 30 seeds, and points are sampled
uniformly without replacement.
\emph{Multiple-survey scenarios:}
nine deterministic, structured layouts model survey unions with overlapping or touching
rectangular sub-surveys, each with its own stride.
This family is motivated by the common structure of real seismic
acquisitions: 5D seismic data (inline, crossline, offset, azimuth, time)
typically decomposes into a (irregular) 2D collection of 3D
sub-surveys, each covering a contiguous patch of the inline--crossline
plane with its own sampling geometry.
They are synthetic but practically motivated stress cases rather than
direct captures from production surveys.
The first six scenarios are effectively 2D (singleton third dimension),
whereas the last three are 3D block arrangements.
After deduplication, they contain between 71\,100 and 9\,264\,024 unique
points.
A separately reported, opt-in fifth-dimensional stress scenario
(\texttt{multiple-survey-9-5d}) extends the same family to genuine
$d=5$ seismic-style geometries at $\sim\!9.2\times 10^7$ unique points;
it is described and measured in
Section~\ref{section:results-stress-5d}.

\begin{table}[ht]
\centering
\caption{Structured multiple-survey scenarios.}
\Description{A table listing the nine structured multiple-survey scenarios, including their short label, number of sub-surveys, effective dimension, unique point count after deduplication, and a short description of the geometric pattern.}
\label{table:structured-scenarios}
\small
\begin{tabular}{lrrrp{7cm}}
\toprule
Label & Surveys & Dim. & Points & Pattern \\
\midrule
2-L
  & 2 & 2D & 71\,100
  & two orthogonal sub-surveys forming an L-shape \\
3-steps
  & 3 & 2D & 192\,200
  & three staggered blocks with stepped offsets \\
4-overlap
  & 4 & 2D & 582\,208
  & four partially overlapping survey blocks \\
5-mixed
  & 5 & 2D & 869\,476
  & mixed block sizes and mixed strides \\
6-bands
  & 6 & 2D & 948\,282
  & six band-like strips with mild overlap \\
7-large
  & 7 & 2D & 1\,308\,196
  & seven larger blocks with one central overlap region \\
12-3D
  & 12 & 3D & 87\,880
  & $3\times4$ 3D block layout \\
16-3D
  & 16 & 3D & 964\,500
  & $4\times4$ 3D block layout \\
36-3D
  & 36 & 3D & 9\,264\,024
  & $6\times6$ 3D block layout \\
\bottomrule
\end{tabular}
\end{table}

\subsection{Metrics}
\label{section:metrics}

For each scenario, configuration, and platform the benchmark records:
construction time (ns/point), structure footprint (bytes), cuboid count
($\kappa$ for IPTs), and four query timings.
The \texttt{pos\_random} and \texttt{at\_random} measurements use
$\min(10^3, n)$ uniformly sampled points or indices.
The \texttt{pos\_all} and \texttt{at\_all} measurements use evenly spaced
samples over the full point or index range, capped at $10^5$ queries for
the standard suite and at $10^4$ for the dense occupancy baselines on the
heavy structured cases.
Each reported per-operation timing is the wall-clock duration of a tight
loop over the entire sample, divided by the sample size; this amortises
the cold-cache and branch-predictor cost of the first iterations into the
steady-state measurement, and the loop itself acts as an implicit warm-up
of the i-cache and predictor state.
For samples below 100 operations--where this implicit warm-up is
insufficient--an explicit warm-up pass over the same sample precedes the
timed pass.
Returned values are consumed via XOR into a sink variable to prevent
dead-code elimination.
Each reported (scenario, configuration, platform) cell is run with 30
independent seeds. For the grid scenarios these seeds govern both the
random point-subset selection and the query sample; for the deterministic
multiple-survey scenarios the underlying point geometry is fixed and the
seeds govern only the random query sampling.
The plots in Section~\ref{section:results} report the arithmetic (or
geometric, where indicated) mean across seeds, with error bars or
shaded bands giving the 10th and 90th percentiles of the underlying
sample distribution to expose the dispersion that scheduler jitter,
allocator placement, and sample-set composition contribute to each
measurement.
We do not pin threads, force NUMA placement, or configure huge pages.
The measurements therefore target robust end-to-end throughput
comparisons under ordinary server conditions rather than cycle-accurate
microarchitectural microbenchmarks.

\paragraph{Cache configurations.}
The implementation exposes four boolean compile-time flags
\texttt{c}~(cuboid size), \texttt{r}~(ruler size), \texttt{e}~(entry's
end index, i.e.\ the prefix-sum upper bound), and \texttt{l}~(entry's
least upper bound, i.e.\ the cached \(\operatorname{lub}\) point used
by \texttt{pos}'s binary search), each controlling whether the
corresponding derived quantity is materialised as an additional stored
field, yielding $2^4 = 16$ possible configurations. We refer to a
configuration by the bit pattern of these flags, e.g.\ \texttt{c0r1e1l0}
for the default with ruler size and entry-end cached but cuboid size
and entry lub omitted.

\paragraph{Persistence and mmap measurements.}
A dedicated storage benchmark exercises the flat-layout requirement of
Section~\ref{section:requirements}. For every measured layout the
benchmark performs the following sequence: (i) build the structure
from the input points; (ii) serialize it to a single contiguous file
via the layout's \texttt{write} primitive; (iii) record the resulting
bytes-on-disk; (iv) drop the operating system's page cache for that
file (\texttt{posix\_fadvise(POSIX\_FADV\_DONTNEED)}, a no-op on
\texttt{tmpfs} but effective on regular file systems); (v) reopen the
file and \texttt{mmap} it read-only into the address space, recording
the load duration; (vi) issue one \texttt{pos}, one \texttt{at}, and
one \texttt{select} call against the freshly mapped instance, with no
prior warm-up, to capture the cold first-response latency that a
live-loaded process would observe; (vii) issue the same
full \texttt{pos}/\texttt{at} sample loops and \texttt{select} sweep
used in the regular benchmark, but executed through the mmap-backed
instance, so that steady-state throughput on the file-backed layout is
directly comparable to the in-memory numbers reported in
Sections~\ref{section:results-pos}--\ref{section:results-selection}.
No pointer fix-up or deserialization step is performed between
steps~(v) and~(vi): the layout is consumed in place from the mapped
bytes, which is the property the IPT shares with the
\textsc{SortedPoints}, \textsc{DenseBitset}, \textsc{BlockBitmap},
\textsc{Roaring}, \textsc{LexRun} and \textsc{RowCSR-k\(*\)} baselines.
This sequence is repeated for three independent random query seeds on
each of the seven smaller structured multiple-survey scenarios
(\texttt{2-l} through \texttt{8-threed}); the larger scenarios are
excluded because the dense bitmap baselines either exceed available
memory or do not finish within the per-scenario budget at full grid
volume. The storage benchmark is run twice, once at the
minimal-metadata cache configuration \texttt{c0r0e0l0} and once at the
paper default \texttt{c0r1e1l0}, so that the on-disk size advantage of
the minimal IPT is visible alongside the steady-state query advantage
of the cached IPT. Each run produces an aggregate file
(\texttt{storage-c0r0e0l0.txt} and \texttt{storage-c0r1e1l0.txt}
respectively) in the per-platform results directory and is invoked via
\texttt{make benchmark-storage}.

\subsection{Configuration Selection}
\label{section:selection}

\paragraph{Construction algorithm.}
On every one of the 61 workload shapes the two construction algorithms
produce identical IPTs (same $\kappa$, same footprint, same entries),
and the construction time agrees within a ${\sim}1.1\times$ band in
either direction across all 20 platform/compiler combinations.
In other words, on non-artificial inputs greedy-plus-merge and
greedy-plus-combine are interchangeable both in result and in
performance: the L-shape rewrite that distinguishes combine never
fires, so it reduces to merge plus a small per-\texttt{add} predicate
check whose remaining wall-clock effect (a few percent in either
direction, with a clean GCC-vs-Clang split) is dominated by codegen
rather than by algorithmic work.
The two algorithms only diverge on the crafted corner-steal family
from Section~\ref{section:construction}, where combine attains the
optimum and merge realises the proven $(d+1)/2$ gap.
On the measured workload suite, combine therefore serves mainly as an
optimality guarantee rather than as a source of measurable throughput
or compression gains.
We therefore use greedy-plus-merge as the simpler default for the
detailed figures below.
Greedy-plus-combine remains important as the optimal-construction variant
proved in Section~\ref{section:construction} and is included in the
artifact and benchmark output. The benchmark binary itself, when
compiled without \texttt{NDEBUG} (i.e.\ as exercised by
\texttt{make verify}), instantiates the crafted corner-steal family
from Section~\ref{section:construction} at startup and verifies that
\texttt{GreedyPlusCombine} recovers the optimum while
\texttt{GreedyPlusMerge} realizes the proven $(d+1)/2$ gap on those
inputs.

\paragraph{Cache configuration.}
We benchmarked all 16 configurations across every platform and scenario.
The entry-end cache (\texttt{e}~flag) remains the dominant
\texttt{at} lever: the cheapest useful variant, \texttt{c0r0e1l0},
reduces geometric-mean \texttt{at} time to $0.485\times$ of the uncached
baseline at only $1.10\times$ memory.
The ruler-size cache (\texttt{r}~flag) is the dominant \texttt{pos}
lever, and the least-upper-bound cache (\texttt{l}~flag), which was
absent from the previous shortlist, contributes additional \texttt{pos}
improvement.
With both \texttt{r} and \texttt{l} enabled, \texttt{pos} drops to about
$0.91\times$ of the uncached baseline, while \texttt{at} stays near
$0.43\times$.
For a query-only objective, \texttt{c0r1e1l1} is therefore the best
frontier point, but its metadata grows to $1.70\times$ the size of the
minimal-footprint \texttt{c0r0e0l0} configuration.
We instead select \texttt{c0r1e1l0} as the default because it keeps the
overhead within a $40\%$ memory budget while retaining most of the query
benefit: relative to \texttt{c0r0e0l0} it yields $0.965\times$
\texttt{pos}, $0.428\times$ \texttt{at}, $1.149\times$ construction
time, and $1.40\times$ metadata.
This makes \texttt{c0r1e1l0} the better paper default when memory is a
first-class constraint rather than a secondary concern.
Particular workloads may still prefer other frontier points, for
example \texttt{c0r1e1l1} on query-heavy deployments with more memory
headroom or \texttt{c0r0e1l0} when minimizing metadata dominates.
No small fixed subset adequately captures the complete trade-off
surface, so we plot all 16 configurations in the selection figure and
use only the selected default in the detailed scenario figures.
Figure~\ref{figure:at-cache} summarises the effect.

\begin{figure}[H]
\centering
\input{generated/at-cache-effect.tex}
\caption{Top: geometric-mean construction, \texttt{pos}, and \texttt{at}
  ratios relative to \texttt{c0r0e0l0}. Bottom: memory
  ratio for the same 16 cache configurations, ordered by increasing
  memory footprint. The paper default \texttt{c0r1e1l0} and the
  query-optimised frontier point \texttt{c0r1e1l1} are shown in bold on
  the x-axis. Top-panel error bars span the 10th to 90th percentile of
  the per-(platform, scenario, seed, grid) ratio distribution.}
\Description{A two-panel cache-configuration summary. The top panel shows grouped bars for construction time, pos time, and at time ratios relative to c0r0e0l0 across all sixteen cache configurations. The bottom panel shows the corresponding memory ratios, with the configurations ordered from lower to higher memory cost. The labels for c0r1e1l0 and c0r1e1l1 are bold to highlight the recommended default and the query-optimized frontier point.}
\label{figure:at-cache}
\end{figure}

The query-side cache rankings are stable across the 20 platform/compiler
combinations: the median Spearman correlation of the 16-configuration
ranking is $\rho = 0.85$ for \texttt{pos} and $\rho = 0.86$ for
\texttt{at}, with minima $0.35$ and $0.54$.
Construction ranks are more platform-sensitive (median $\rho = 0.50$,
minimum $-0.39$), so the detailed construction numbers below are tied to
the reference platform rather than interpreted as a universal platform
ordering.
We use AMD EPYC~9454 (Rocky~10.1 / GCC~14.3.1) as that reference.

\section{Experimental Results}
\label{section:results}

\subsection{Construction Performance}
\label{section:results-construct}

Both the IPT constructors and the \texttt{SortedPoints} baseline are
benchmarked on lexicographically sorted input, matching the precondition
of the optimality proof in Section~\ref{section:construction} and the
scope statement in Section~\ref{section:discussion}; neither side
includes the cost of producing that ordering. Pipelines that start from
an unordered point list incur the same sort cost for both formats, so
the build-time ratio reported here is invariant under that shared
prefix; the absolute amortization point shifts uniformly once a sort is
charged.

On the EPYC~9454 reference platform, greedy-plus-merge construction
reaches $42.2{-}45.8$\,Mpoint/s on regular grids.
On the structured multiple-survey scenarios it achieves
$36.5{-}45.8$\,Mpoint/s, again in a comparatively narrow band despite the
varying block layouts.
\texttt{SortedPoints} construction is essentially a memory copy
and reaches up to ${\sim}\!660$\,Mpoint/s on the regular grids and
$90{-}330$\,Mpoint/s on the structured scenarios, so the IPT is
${\sim}3.8\times$ more expensive to build on the structured workloads
(geometric mean), again with both sides starting from sorted input.
The cost is dominated by cuboid discovery and merging.
At the geometric-mean rates on the structured scenarios
(${\sim}40$\,Mpoint/s construction, ${\sim}10$\,Mpos/s \texttt{pos}),
the $3.8\times$ build overhead amounts to an extra ${\sim}18$\,ns per
input point; the $2.33\times$ \texttt{pos} speedup over
\texttt{SortedPoints} saves ${\sim}125$\,ns per query.
Under the same sorted-input assumption, the IPT therefore recovers its
build cost after approximately $N/7$ \texttt{pos} queries, where $N$ is
the number of input points.
Any workload that issues at least one query per seven stored points
(like repeated-timestep simulations with fixed geometry do) amortises
construction within the first pass; pipelines that must additionally
sort their input pay that sort cost once and amortise it identically
for both formats.
Figure~\ref{figure:construct-multi-survey} shows the per-point
construction time on the nine structured multiple-survey scenarios.
Figure~\ref{figure:construct-vs-density} shows construction time
as a function of density for the four grid sizes.

\begin{figure}[H]
\centering
\input{generated/construct-multiple-survey.tex}
\caption{Construction time per point on the multiple-survey
  scenarios (\texttt{c0r1e1l0}, EPYC~9454 reference platform;
  error bars span the 10th-90th percentile across 30 seeds).}
\Description{A bar chart comparing IPT and SortedPoints construction time per point across the nine structured multiple-survey scenarios on the EPYC 9454 reference platform, with 10th-to-90th-percentile error bars across 30 seeds.}
\label{figure:construct-multi-survey}
\end{figure}

\begin{figure}[H]
\centering
\input{generated/construct-vs-density.tex}
\caption{IPT construction time per point vs.\ point-set density
  for four grid sizes (30-seed mean, \texttt{c0r1e1l0},
  EPYC~9454 platform; error bars span the 10th-90th percentile across
  the 30 seeds).}
\Description{A line plot of IPT construction time per point versus point-set density for four regular grid sizes, with 10th-to-90th-percentile error bars across 30 seeds. The curves show that construction cost decreases and then flattens as density approaches one hundred percent.}
\label{figure:construct-vs-density}
\end{figure}

\subsection{Compression and Memory}
\label{section:results-kappa}

On fully regular grids the IPT produces a single cuboid ($\kappa = 1$)
regardless of grid size.
For random subsets of the $100{\times}100{\times}100$ grid, $\kappa$ is concave in the
point-set density: it rises from ${\sim}7\,100$ at $1\%$ to a peak of
${\sim}195\,000$ at $50\%$ and then declines back to~1 at $100\%$
(30-seed mean, coefficient of variation $< 0.2\%$).
The multiple-survey scenarios yield cuboid counts
ranging from $\sim\!2\,100$ to $\sim\!137\,000$.
The largest count still occurs in the 7-large 2D layout, whereas the
16-block and 36-block 3D layouts compress much better
($\sim\!4\,900$ and $\sim\!3\,000$ cuboids) because their repeated local
block structure exposes long regular runs despite the higher number of
sub-surveys.

Because the IPT stores per-cuboid metadata rather than per-point
data, memory is directly proportional to~$\kappa$.
With $\kappa = 1$ the IPT stores a single entry regardless of the
number of points.
On the application-relevant multiple-survey scenarios the
memory ratio (IPT\,/\,SortedPoints) is $0.138$
for \texttt{c0r1e1l0} and $0.105$ for the minimal
configuration \texttt{c0r0e0l0},
i.e.\ the selected default uses about $14\%$ of the sorted-points
footprint on those structured layouts.
For the sparsest random subsets ($1\%$) the high $\kappa$ can
push the IPT to ${\sim}\!4.0\times$ the sorted-points size;
the minimal-footprint \texttt{c0r0e0l0} reduces that worst case
substantially.

\begin{figure}[H]
\centering
\input{generated/kappa-vs-density.tex}
\caption{Normalized cuboid count $\kappa/N$ as a function of
  point-set density for four grid sizes (30-seed mean, EPYC~9454
  platform; error bars span the 10th-90th percentile across the 30
  seeds).}
\Description{A line plot of the normalized cuboid count kappa over N versus density for four grid sizes, with 10th-to-90th-percentile error bars across 30 seeds. The curves peak at intermediate densities and drop toward zero as the data becomes fully regular.}
\label{figure:kappa}
\end{figure}

\subsection{Query Performance: \texttt{pos}}
\label{section:results-pos}

\emph{IPT vs.\ mixed-radix (regular grids).}
Even without the \texttt{l}-cache, the single-entry binary search on
regular grids remains effectively free at the measured sizes.
On regular grids ($\kappa = 1$) the IPT matches pure mixed-radix
rank/unrank closely: the geometric-mean \texttt{pos} time is only
$0.99\times$ that of the \texttt{Grid} baseline and the geometric-mean
\texttt{at} time is $1.00\times$.
The tighter memory budget therefore does not materially degrade the
regular-grid case.

\emph{IPT vs.\ sorted-points.}
On regular grids the sampled \texttt{pos\_random} benchmark shows the IPT's
\texttt{pos} to be $9.3{-}40.7\times$ faster than \texttt{SortedPoints}
(which performs a full \texttt{lower\_bound} over all $n$ points).
On the regular grids the IPT reaches $169{-}180$\,Mpos/s.
On the multiple-survey scenarios the advantage is
${\sim}2.33\times$ (geometric mean), with the IPT reaching
$5.0{-}22.6$\,Mpos/s.
As random-subset density drops, $\kappa$ rises and the IPT's
advantage shrinks: at $1\%$ and $50\%$ density the IPT is still
slower on the $100{\times}100{\times}100$ grid (ratios $1.22$ and
$1.45$), the crossover occurs between $75\%$ and $90\%$ density
(ratios $1.17$ and $0.86$), and by $99\%$ density the IPT is already
about $1.87\times$ faster owing to the drop in~$\kappa$.
The exhaustive \texttt{pos\_all} timings show the same ranking and even
stronger separation on the structured scenarios, where the IPT is
more than $2\times$ faster than \texttt{SortedPoints} in geometric mean.
Figure~\ref{figure:pos-multi-survey} shows the absolute \texttt{pos}
performance on the nine structured multiple-survey scenarios.
Figure~\ref{figure:pos-ratio} shows the full density sweep.

Among the additional structured baselines, \texttt{RowCSR-k1} is the
only cache-independent layout that improves \texttt{at} over the IPT,
but it still loses on \texttt{pos} and uses substantially more memory.
\texttt{RowCSR-k2} and \texttt{LexRun} also remain slower than the IPT
on both queries across the structured scenarios.
The dense occupancy baselines are only tractable on the seven smaller
structured scenarios, and there they remain two to three orders of
magnitude slower on both queries despite bit-level storage.

\begin{figure}[H]
\centering
\input{generated/pos-multiple-survey.tex}
\caption{Absolute \texttt{pos} performance on the multiple-survey
  scenarios (\texttt{c0r1e1l0}, EPYC~9454 reference platform;
  error bars span the 10th-90th percentile across 30 seeds).}
\Description{A bar chart comparing absolute pos query performance for IPT and SortedPoints across the nine structured multiple-survey scenarios on the EPYC 9454 reference platform, with 10th-to-90th-percentile error bars across 30 seeds.}
\label{figure:pos-multi-survey}
\end{figure}

\begin{figure}[H]
\centering
\input{generated/pos-ratio-vs-density.tex}
\caption{\texttt{pos} time ratio (IPT / SortedPoints) vs.\ density
  (\texttt{c0r1e1l0}, EPYC~9454 platform, 30-seed geometric mean;
  error bars span the 10th-90th percentile across the 30 seeds).}
\Description{A line plot of the pos time ratio between IPT and SortedPoints versus point-set density for four grid sizes, with 10th-to-90th-percentile error bars across 30 seeds. The ratio is above one at low and medium densities and drops well below one near full regularity.}
\label{figure:pos-ratio}
\end{figure}

\subsection{Query Performance: \texttt{at}}
\label{section:results-at}

\texttt{SortedPoints} answers \texttt{at} by direct array indexing
($O(1)$). The IPT requires a binary search plus Horner unrank.
On regular grids the IPT is $1.05{-}8.69\times$ slower than
\texttt{SortedPoints}, but it stays within ${\sim}1\%$ of the ideal
regular-grid baseline because the selected caches leave the one-entry
search near-trivial.
On random subsets the gap widens to roughly $25{-}45\times$ because $\kappa$
determines the binary-search depth.
This is the expected cost of trading per-point storage for
per-cuboid metadata.
Figure~\ref{figure:at-multi-survey} shows the absolute \texttt{at}
performance on the nine structured multiple-survey scenarios.
On those structured scenarios the IPT reaches $8.1{-}28.3$\,M\texttt{at}/s
and is $18.7\times$ slower than \texttt{SortedPoints} in geometric mean.
The exhaustive \texttt{at\_all} timings show the same ranking but a
smaller gap, because the evenly spaced index pattern has
better locality for both layouts.

\begin{figure}[H]
\centering
\input{generated/at-multiple-survey.tex}
\caption{Absolute \texttt{at} performance on the multiple-survey
  scenarios (\texttt{c0r1e1l0}, EPYC~9454 reference platform;
  error bars span the 10th-90th percentile across 30 seeds).}
\Description{A bar chart comparing absolute at query performance for IPT and SortedPoints across the nine structured multiple-survey scenarios on the EPYC 9454 reference platform, with 10th-to-90th-percentile error bars across 30 seeds.}
\label{figure:at-multi-survey}
\end{figure}

\subsection{Selection and Restriction}
\label{section:results-selection}

The cuboid algebra of Section~\ref{section:operations-restriction}
gives the IPT a closed-form selection operator
$\textsc{select}: Q \mapsto \{\, C_h \cap Q : C_h \cap Q \neq \emptyset
\,\}$
and, by re-feeding its result into the from-range constructor, a
restriction operator $\textsc{restrict}(Q) = \mathrm{IPT}(\textsc{select}(Q))$.
Under the same tuple-comparison model used in
Section~\ref{section:operations-restriction}, the expected cost of one
full enumeration is
\[
  O(d\log\kappa + d\,t + d\,r),
\]
where $t$ is the number of cuboid candidates touched by the binary-bound
search and $r$ is the number of non-empty intersection cuboids
returned. The $d\log\kappa$ term is the one-shot binary search, charging
$O(d)$ per lexicographic tuple comparison; $d\,t$ is the per-candidate
intersection cost; $d\,r$ is the per-result yield cost. This matches
the $O(d\log\kappa + d\,t)$ bound stated in
Section~\ref{section:operations-restriction} (with $r\le t$ absorbed
into the $d\,t$ term) and additionally exposes the yield term that the
select benchmark below isolates.

To validate this decomposition we instrumented the same seven baselines
used for \texttt{at}/\texttt{pos} (\textsc{Grid}, \textsc{SortedPoints},
\textsc{DenseBitset}, \textsc{BlockBitmap}, \textsc{LexRun},
\textsc{RowCSR-k1}, \textsc{RowCSR-k2}) and the three IPT
configurations (\textsc{Regular}, \textsc{GreedyPlusCombine},
\textsc{GreedyPlusMerge}) with a uniform \texttt{select} method that
emits a lazy stream of \texttt{Cuboid<D>} fragments
covering exactly the points in the query window. The point-level
baselines emit one singleton cuboid per matching point; \textsc{LexRun}
emits one cuboid per surviving lex-run; \textsc{Grid} and the IPT
variants emit at most one fragment per intrinsic cuboid. The benchmark
binary, when compiled without \texttt{NDEBUG} (i.e.\ as exercised by
\texttt{make verify}), additionally runs $200$ random axis-aligned
queries on a $16\!\times\!16\!\times\!16$ grid at startup and checks
that all ten implementations yield identical point sets after
enumeration.

\paragraph{What the \texttt{restrict} numbers actually compare.}
The \texttt{restrict} pipeline reported here is, for every
implementation including the non-IPT baselines, the same two-stage
recipe \emph{select-then-IPT-rebuild}: the implementation's
\texttt{select} stream is fed into the IPT from-range constructor and
the resulting sub-IPT is the measured output. This is intentional, and
matches the IPT's native restriction operator
$\textsc{restrict}(Q) = \mathrm{IPT}(\textsc{select}(Q))$ from
Section~\ref{section:operations-restriction}; for the non-IPT baselines
it is \emph{not} a measurement of any layout-native restriction
operator. \textsc{SortedPoints}, \textsc{LexRun}, \textsc{RowCSR-k\(*\)},
\textsc{DenseBitset} and \textsc{BlockBitmap} are not closed under
their own restriction in this representation: they would naturally
produce a smaller instance of themselves, not an IPT. The numbers
therefore quantify ``how fast can each layout's selection feed an IPT
sub-domain rebuild'', which is the shared workload the IPT advertises;
they should not be read as a comparison of layout-native restriction
operators.

The full 20-platform sweep emits \texttt{select\_d*} and
\texttt{restrict\_d*} measurement lines (for $d \in \{0.001, 0.01, 0.1,
1.0\}$ relative window volumes) for every algorithm on every scenario,
so the asymptotic claims above are checked across the same matrix used
for \texttt{at} and \texttt{pos}.
The measured \texttt{select} time scales linearly with the number of
yielded fragments for every implementation, in line with the $O(d\,r)$
yield term.
On the EPYC~9454 reference platform the \texttt{restrict} cost (one full
\texttt{select} pipe followed by a from-range IPT construction) is
only about $1.6\%$ above the bare \texttt{select} cost for the IPT in
geometric mean across the structured-scenario density sweep,
confirming that restriction does not introduce
asymptotic overhead beyond enumeration.

The same comparison against the other instrumented baselines confirms
the ranking established for \texttt{pos}.
At a $d = 0.01$ window on the nine structured multiple-survey
scenarios, the IPT's \texttt{select} is geometric-mean
$7.0\times$ faster than \texttt{SortedPoints},
$7.6\times$ faster than \texttt{RowCSR-k1},
$68\times$ faster than \texttt{RowCSR-k2}, and
$215\times$ faster than \texttt{LexRun};
on the seven smaller scenarios where they are tractable,
\texttt{DenseBitset} and \texttt{BlockBitmap} are
$10.9\times$ and $8.8\times$ slower than the IPT.
The corresponding \texttt{restrict} ratios fall within $\pm 15\%$ of
those values (e.g.\ $6.2\times$ for \texttt{SortedPoints} and
$6.7\times$ for \texttt{RowCSR-k1}), since restriction adds only the
constant-factor from-range rebuild on top of \texttt{select}.
The gaps widen further at smaller $d$ (where the IPT's binary search
and skip over empty cuboids dominate the entire pipeline) and shrink
at $d = 1$ as every implementation becomes yield-bound.

\subsection{Persistence: Serialization, mmap-Load, and First-Response Latency}
\label{section:results-storage}

This subsection reports the measurements collected by the persistence
benchmark of Section~\ref{section:metrics}. The goal is to back the
flat-layout requirement of Section~\ref{section:requirements} with
empirical evidence: that the IPT and the comparable baselines can be
serialized to a single contiguous file, mapped back into the address
space, and queried directly from the mapped bytes with no
reconstruction step.

\paragraph{What is measured.} For each of the seven smaller structured
scenarios (\texttt{multiple-survey-2-l}, \texttt{3-steps},
\texttt{4-overlap}, \texttt{5-mixed}, \texttt{6-bands},
\texttt{7-large}, \texttt{8-threed}) and each of the eight
storage-capable layouts (\textsc{IPT} via greedy-plus-merge,
\textsc{SortedPoints}, \textsc{DenseBitset}, \textsc{BlockBitmap},
\textsc{Roaring}, \textsc{LexRun}, \textsc{RowCSR-k1},
\textsc{RowCSR-k\(d{-}1\)}), the benchmark emits five derived
quantities per (cache configuration, seed) cell:
(i)~\emph{serialize}, the wall-clock time to write the structure to
disk in nanoseconds per stored point; (ii)~\emph{bytes-on-disk}, the
final file size; (iii)~\emph{cold-load}, the wall-clock time to drop
the page cache, reopen the file, and \texttt{mmap} it; (iv)~the
first-response latency of \texttt{pos}, \texttt{at}, and
\texttt{select} on the freshly mapped instance, with no warm-up;
(v)~the steady-state \texttt{pos}, \texttt{at}, and \texttt{select}
throughput measured through the mmap-backed instance using the same
sample loops as Sections~\ref{section:results-pos}--%
\ref{section:results-selection}.

\paragraph{What the experiment is designed to show.} Three properties
are under test. First, that the on-disk size of the IPT tracks the
in-memory metadata footprint reported in
Section~\ref{section:results-summary}, and in particular that the
minimal-cache \texttt{c0r0e0l0} configuration is the smallest of the
eight measured layouts on the structured-regime workloads where it
applies. Second, that the cold first-response latency of the IPT after
an \texttt{mmap} is bounded by the page-fault and binary-search cost
of a single query, and stays within roughly an order of magnitude of
the steady-state \texttt{pos}/\texttt{at} numbers reported in the
in-memory subsections, rather than being dominated by a
reconstruction or pointer-fix-up phase. Third, that the steady-state
\texttt{pos}/\texttt{at}/\texttt{select} throughput on the
file-backed instance is comparable (within typical OS-level paging
noise) to the in-memory numbers, confirming that mmap is not just a
load mechanism but also a viable steady-state runtime substrate for
the IPT. The \texttt{c0r0e0l0} versus \texttt{c0r1e1l0} pairing
additionally exposes the size/throughput trade-off discussed in
Section~\ref{section:selection} from the persistence side: the same
default that is preferable in memory is also the preferable target on
disk, because the cache fields are stored once and amortised over all
future cold loads of the same file.

The per-scenario, per-algorithm tables and the corresponding figure
are produced by a dedicated extractor (\texttt{plot/storage-summary})
over the aggregate files \texttt{storage-c0r0e0l0.txt} and
\texttt{storage-c0r1e1l0.txt}; both files are emitted into each
platform's results directory by \texttt{make benchmark-storage}.

\subsection{Hit-vs-Miss \texttt{pos} Cost}
\label{section:results-pos-mix}

The throughput of \texttt{pos} reported throughout
Section~\ref{section:results} measures \emph{hits} only: every probed
point is in the index, every algorithm returns a valid linear
position, and no exception is ever thrown. Real call sites are not
that pure; even a small fraction of misses can dominate the cost if a
miss is signalled by an exception, because exception unwinding on
mainstream toolchains costs several thousand cycles and dwarfs a
plain \texttt{pos} call. To avoid blending an unwind cost into the
hit-only timings the artifact exposes a non-throwing parallel API
\texttt{try\_pos} that returns \texttt{std::optional<Index>} and is
implemented for every algorithm without changing the existing,
throwing \texttt{pos}; both APIs share the same internal lookup so
their hit-only ns/op are within measurement noise of each other.

The hit-vs-miss benchmark (\texttt{make benchmark-pos-mix}) calls
\texttt{try\_pos} on a deterministic mixed sample of size
$1\,000$ for each of five miss ratios
$\{0.0,\,0.01,\,0.10,\,0.50,\,1.00\}$ and for two miss kinds: an
\emph{in-grid} miss (a coordinate that lands at a valid grid node but
is not in the stored set, generated by rejection-sampling the lattice
$origin + i\cdot stride$) and an \emph{out-of-grid} miss (one
coordinate perturbed to $lo - stride$ or $lo + extents\cdot stride$).
The seven structured 3D scenarios used by the storage benchmark
(\texttt{multiple-survey-2-l} through \texttt{multiple-survey-8-threed})
are run for three seeds each. Aggregates per algorithm and per miss
ratio are the geometric mean across the seven scenarios after
averaging over seeds; the resulting curves on the reference platform
are shown in Figure~\ref{figure:pos-mix-summary}.

\begin{figure}[ht]
\centering
\input{generated/pos-mix-summary}
\Description{Two-panel figure of average try_pos cost in nanoseconds
  per query as a function of miss ratio for in-grid and out-of-grid
  misses, one curve per algorithm, on the reference platform.}
\caption{Hit-vs-miss \texttt{try\_pos} cost (geometric mean over the
  seven structured 3D scenarios, averaged across three seeds, on the
  reference platform). Left panel: misses that land at a valid grid
  node. Right panel: misses with one coordinate outside the bounding
  grid. Algorithms whose miss path performs a popcount over a dense
  bitset (DenseBitset, BlockBitmap) carry the full structure cost on
  every hit but become very cheap as the miss ratio approaches one;
  IPT, SortedPoints, Roaring, LexRun and RowCSR all show
  near-constant cost across the sweep.}
\label{figure:pos-mix-summary}
\end{figure}

Three observations follow from the data. First, IPT's
\texttt{try\_pos} cost is essentially flat in the miss ratio: the
ruler chain settles a miss at the same lower-bound step it would use
for a hit, and a missed entry returns \texttt{std::nullopt} without
descending into the dense block. Second, the popcount-based bitmap
algorithms (DenseBitset and BlockBitmap) pay their full structure
cost on every hit but become very cheap as the miss ratio approaches
one, because the bit-test that fails the membership check short-cuts
the popcount; this exposes a workload-dependent crossover that the
hit-only \texttt{pos} number could not show. Third, all CSR-style and
runlength-style algorithms (SortedPoints, Roaring, LexRun, RowCSR)
see only a small change in per-call cost across the sweep, because
their miss path is a binary search that fails one comparison earlier
than a hit. The two miss kinds (in-grid vs.\ out-of-grid) differ by
less than a factor of two for every algorithm; out-of-grid misses are
slightly cheaper because they fail the per-axis bounds check before
the membership query is built.

\subsection{Higher-Dimensional Stress: a 5D Multiple-Survey Scenario}
\label{section:results-stress-5d}

To check that the IPT properties survive at higher~$d$ and at a point
count comparable to a production seismic geometry, the artifact
includes a fifth, opt-in stress scenario \texttt{multiple-survey-9-5d}
(executed by \texttt{make benchmark-5D}). It places a $3\times 4$
arrangement of twelve 3D sub-surveys, each carried by its own strides
in the inline/crossline/offset axes and each expanded by a fixed
$125\times 80$ azimuth--time block, into a bounding grid of
$47\times 44\times 41\times 125\times 80$. After deduplication of the
overlap between neighboring sub-surveys this yields $92\,260\,000$
unique 5D points. We collected this scenario on the same 18 of the 20
platform/compiler combinations that supplied a stress run with the
default \texttt{c0r1e1l0} cache configuration. The aggregate appears
in Table~\ref{table:stress-5d-summary}.

\begin{table}[ht]
\centering
\caption{Aggregate \texttt{multiple-survey-9-5d} measurements
  (greedy-plus-merge, \texttt{c0r1e1l0}, geometric mean across the 18
  stress platforms; bracketed range is the 10th--90th percentile).
  Cuboid count and IPT byte size are deterministic functions of the
  input and therefore platform-independent.}
\Description{Table of aggregate measurements for the
  multiple-survey-9-5d 5D stress scenario, summarising construction
  rate, pos and at throughput, and select/restrict latency at two
  density levels across 18 platform/compiler combinations.}
\label{table:stress-5d-summary}
\input{generated/stress-summary.tex}
\end{table}

Three observations carry over from the lower-dimensional results.
First, the cuboid count is $\kappa = 2154$, so the IPT metadata fits in
$379\,104$ bytes ($\approx 370\,$KiB) for $9.2\times 10^7$ points; the
IPT therefore stores about $4\,300$ points per cuboid record on this
geometry, well below $N$ and consistent with the structured regime
identified in Section~\ref{section:results-summary}. Second, the
\texttt{at}/\texttt{pos} gap (\texttt{at\_random} faster than
\texttt{pos\_random}) persists at $d=5$, in line with the per-query
analysis. Third, \texttt{select} and \texttt{restrict} stay in the
$10$--$10^3\,\mu$s range up to a $d=0.1$ relative window even at this
size, confirming that the cuboid-algebra path scales with the structure
of the data rather than with $N$.

The IPT-only stress configuration deliberately omits the
non-IPT baselines: at this size and dimensionality, several of them
exceed available memory or do not finish within the per-scenario
budget. The numbers above therefore quantify the IPT's behavior at
$d=5$ and are not a head-to-head comparison; they should be read as
evidence that the structured-regime properties demonstrated in
Sections~\ref{section:results-pos}--\ref{section:results-selection}
do not collapse when $d$ and $N$ are pushed simultaneously.

\subsection{Summary of Empirical Regimes}
\label{section:results-summary}

The results show three practical regimes. On (strided-)regular inputs,
the IPT preserves the regular-grid structure and, in our benchmarks,
nearly matches the ideal mixed-radix baselines while keeping metadata
constant in~$d$.
On structured irregular inputs such as the multiple-survey workloads,
the IPT keeps $\kappa$ far below $N$, which leads to substantial memory
savings and clear \texttt{pos} gains over \texttt{SortedPoints}.
On highly fragmented random subsets, $\kappa$ approaches~$N$, the memory
advantage disappears, and the cost of the global binary search dominates
\texttt{at} and can also erase the \texttt{pos} advantage. The governing
parameter across all scenarios is therefore not sparsity alone but the
amount of geometric structure captured by a small cuboid count.

\section{Discussion and Limitations}
\label{section:discussion}

The experimental results are consistent with the complexity analysis.
Construction scales linearly in~$N$
(Section~\ref{section:results-construct}).
The cuboid count $\kappa$ grows sub-linearly on random subsets and
collapses to~1 on (strided-)regular inputs
(Section~\ref{section:results-kappa}). Query times show the expected
dependence on~$\kappa$ for both \texttt{at} and \texttt{pos}
(Sections~\ref{section:results-pos} and~\ref{section:results-at}).

Across all 61 workload shapes, the two construction algorithms produced
identical decompositions and construction times that agree within a
${\sim}1.1\times$ band in either direction across the 20
platform/compiler combinations. This does not contradict the
counterexamples from Section~\ref{section:construction}: it confirms
that the corner-steal patterns required for combine to do anything
different from merge are rare enough to be effectively absent from the
tested input family. On non-artificial inputs the two algorithms are
interchangeable in both result and performance, so we keep
greedy-plus-merge as the simpler practical default for the detailed
figures and rely on greedy-plus-combine only as the
optimal-construction variant verified by \texttt{make verify}. On this
benchmark suite its practical role is therefore correctness insurance
and regression checking, not additional measured speed or compression.

The evaluation also clarifies the paper's effective regime. The method
is most compelling when irregular geometry still decomposes into few
cuboids. In that setting it combines exactness, compact metadata, and
fast \texttt{pos}. When fragmentation drives $\kappa$ toward~$N$, the
advantages diminish and the representation approaches the cost profile
of explicit point storage.

A favorable but only partially covered regime in our
benchmarks is a
\emph{common degenerate dimension}. If all cuboids share the same
ruler in one coordinate direction, then enlarging that common ruler
does not change~$\kappa$. The IPT therefore keeps the same number of
cuboids, the same metadata footprint, and the same asymptotic
\texttt{at} and \texttt{pos} costs. The multiple-survey benchmarks do
not isolate this effect cleanly: the first six structured scenarios use a
singleton third dimension and are effectively two-dimensional, whereas
the last three scenarios are 3D and vary in every dimension. If the
common third-dimension extent in the first family were increased from
$1$ to $1000$, the IPT metadata would remain unchanged, whereas a
\texttt{SortedPoints} representation would store $1000$ times as many
points and its \texttt{pos} search would become more expensive because it
scales with the total point count.

The main limitation is fragmentation. As $\kappa$ approaches~$N$, the
advantages diminish. The flat IPT is also not a nearest-neighbor
structure, and dynamic updates are awkward in the frozen form because
they induce nonlocal changes in prefix indices.

One scope limit bears stating explicitly: the
construction algorithms are analyzed for sorted input and are evaluated
under that assumption throughout the experiments.

The experimental scenarios cover regular grids, random subsets,
and structured synthetic survey layouts, including both 2D-in-3D
multiple-survey cases and genuinely 3D block-survey layouts. That
range is broad enough to
show the main trade-offs, but it does not replace application studies
on production geometries. Real-world workloads may expose additional
construction patterns, especially outside the geoscience setting.

\paragraph{Dynamic and AMR settings.}
The flat IPT is constructed once from sorted input and is read-only
afterward. In dynamic or adaptive-mesh-refinement (AMR) settings where
the topology changes frequently, the relevant cost is not just one
construction but the amortized cost of repeated re-builds. The linear
$O(dN)$ construction (Section~\ref{section:results-construct}) keeps a
full re-build practical when topology changes happen on a much coarser
time scale than queries, but it is not by itself a streaming or
incremental update procedure. A useful incremental variant would have
to amortize prefix-sum recomputation and cuboid-list editing across
small local edits; this is left to future work and is the most
prominent gap between the present design and a fully dynamic spatial
index.

\paragraph{Practical $d$-scaling.}
The asymptotic costs $O(d+\log\kappa)$ for \texttt{at} and
$O(d\log\kappa)$ for \texttt{pos} are linear in~$d$, but the
constant-factor picture changes as $d$ grows beyond the dimensions
covered here ($d\le 3$). Each cuboid record stores $O(d)$ ruler
fields, so for fixed~$\kappa$ the metadata footprint per cuboid grows
linearly with~$d$ and an entire cuboid record stops fitting in a
single cache line at modest $d$ (roughly $d\gtrsim 4$ on the platforms
studied here, depending on cache configuration). Beyond that point the
binary search over $\operatorname{lub}(C_h)$ pays an additional cache
miss per comparison, and the practical \texttt{pos} cost is closer to
$O(d\log\kappa)$ memory accesses than to $O(d\log\kappa)$ register
operations. The structural advantages of the flat IPT (exact geometry,
geometric restriction, pointer-free layout) all remain, but the
benchmarked throughput numbers should not be extrapolated unchanged to
high-dimensional settings.

\subsection{Reproducibility}
\label{section:reproducibility}

The full benchmark and paper-build pipeline is included with the
submission and is also published as an open-source artifact under the
license stated in the \texttt{LICENSE} file. The top-level
\texttt{Makefile} drives the entire pipeline:
\begin{itemize}
  \item \texttt{make verify} compiles all 16 cache-flag combinations
    with assertions enabled and runs the corner-steal and
    cross-implementation selection verifiers on a reduced scenario
    subset.
  \item \texttt{make benchmark} runs the cache-dependent IPT
    implementations under all 16 configurations and the
    cache-independent baselines once per platform; output goes to
    \texttt{results/<platform>/}.
  \item \texttt{make benchmark-stress} runs the
    \texttt{multiple-survey-9-5d} 5D stress scenario at the
    \texttt{c0r1e1l0} default
    (Section~\ref{section:results-stress-5d}).
  \item \texttt{make plots} regenerates every figure in this paper
    by compiling the per-figure extractors in \texttt{plot/*.cpp}, running
    them over \texttt{results/}, and invoking gnuplot on the
    \texttt{plot/*.gp} scripts.
  \item \texttt{make paper} re-runs \texttt{pdflatex} and
    \texttt{bibtex} on \texttt{IPT.tex}.
\end{itemize}
The platforms used in this paper (Table~\ref{table:platforms}) are
listed by OS, kernel, CPU, and compiler; each result directory under
\texttt{results/} encodes the same identifying string so that an
independent re-run on the same hardware/compiler combination produces
a directly comparable subdirectory. Compilers are pinned by version in
those directory names. The header-only IPT library lives in
\texttt{ipt/}; the benchmark driver is split across
\texttt{benchmark/Common.hpp}, \texttt{benchmark/Scenarios.hpp}, and
one translation unit per scenario under
\texttt{benchmark/scenarios/}; the per-figure extractors and gnuplot
scripts are colocated in \texttt{plot/}; aggregated tabular outputs
land in \texttt{generated/}. Citation metadata for the artifact is
provided in \texttt{CITATION.cff}. The single external dependency
required by the benchmark driver is the CRoaring amalgamation
(v4.6.1, dual Apache-2.0\,/\,MIT licensed), vendored unmodified under
\path{third_party/CRoaring/} and built once into
\path{build/benchmark/libroaring.a} that every scenario binary links
against; no system installation is required.

\section{Conclusion}
\label{section:conclusion}

We introduced the Adaptive Index--Point Translator (IPT), an exact,
adaptive, flat representation for irregular spatial point sets that
composes a cuboid-selection binary search with the classical Horner
rank/unrank scheme. The representation provides exact geometry,
adaptive storage, flat memory layout, and geometric restriction within
one array-based structure.

On the algorithmic side, the paper gives linear-time construction
procedures based on greedy repair. In particular, greedy-plus-combine
achieves the minimum cuboid count in $O(dN)$ time on sorted input. The
design-space analysis explains why the flat IPT is a natural
systems-oriented extension of Horner-style rank/unrank for the goals of
this paper.

The experiments show that the approach is most effective when
irregular data retains substantial cuboid structure and when the
dominant query is \texttt{pos}-heavy, selection-heavy, or
restriction-heavy. In that regime the
IPT keeps about $14\%$ of the sorted-points memory on the tested
structured survey layouts, improves \texttt{pos} by about $2.33\times$
on those layouts, and offers a closed-form geometric restriction
operator. It also preserves the regular-grid case closely with
the \texttt{c0r1e1l0} cache configuration. On \texttt{at}-heavy
traversals of the same structured layouts the IPT is slower than a
sorted-points table; the experiments quantify this gap rather than
claim a uniform speedup. Across the 61 organic workload shapes the
proven-optimal greedy-plus-combine constructor coincided exactly with
the simpler greedy-plus-merge constructor; its practical role in this
benchmark suite is therefore correctness insurance for the corner-steal
family, not additional measured speed or compression.
The cache-sweep results confirm that \texttt{c0r1e1l1} achieves the
best pure query performance, but \texttt{c0r1e1l0} is the more broadly
usable default because it caps metadata overhead at $40\%$ while
retaining most of that query benefit. Particular workloads can still
justify other cache settings, and the full 16-configuration sweep
characterises the complete trade-off space.
Future work includes broader application studies on production point
distributions outside the geoscience domain, broader empirical
comparisons against additional exact index structures, and dynamic or
incremental construction strategies for mutable data.

\appendix

\section{Proofs}

\subsection{One-Cuboid Recovery}
\label{proof:one-cuboid-recovery}

We prove that if the input point set is representable by a single cuboid and is processed in lexicographic order, the greedy-plus-merge algorithm reconstructs a one-cuboid decomposition with at most~$d$ intermediate cuboids at any point.

\begin{proof}
By induction on~$d$.

\emph{Base case ($d=1$).}
A one-dimensional cuboid is a ruler $R=[b,s,e)=\{b,b+s,\ldots,b+s(n-1)\}$. Its points arrive in lexicographic order exactly in that order. After the first $t$ points have been processed, the greedy-plus-merge state is the single ruler $[b,s,b+st)$. The next point $b+st$ extends this ruler to $[b,s,b+s(t+1))$, so the algorithm never creates a second cuboid. Hence it finishes with one cuboid, and at every moment there is at most one intermediate cuboid.

\emph{Inductive step ($d \geq 2$).}
A $d$-dimensional cuboid $C = R_0 \times \cdots \times R_{d-1}$ with $R_0=[b_0,s_0,e_0)$ decomposes in lexicographic order into the consecutive slices
\[
  S_h = \{b_0 + s_0 h\} \times R_1 \times \cdots \times R_{d-1},
  \qquad h=0,\ldots,|R_0|-1.
\]
For each $h$, let
\[
  R_0^{(h)} = [b_0,s_0,b_0+s_0(h+1)),
  \qquad
  A_h = R_0^{(h)} \times R_1 \times \cdots \times R_{d-1}.
\]
Thus $A_h$ is the cuboid formed by the first $h+1$ slices.

We claim by induction on $h$ that after the algorithm finishes slice $S_h$, it has reconstructed exactly the cuboid $A_h$. We also claim that while slice $S_h$ is being processed, the cuboid list consists of at most one accumulated cuboid from earlier slices together with the cuboids currently created inside $S_h$. Hence the total number of cuboids present at that time is at most $1+(d{-}1)=d$.

For $h=0$, the slice $S_0$ is, apart from its fixed first coordinate, a $(d{-}1)$-dimensional one-cuboid input $R_1 \times \cdots \times R_{d-1}$. By the induction hypothesis on~$d$, the greedy-plus-merge algorithm reconstructs $S_0$ as one cuboid while using at most $d{-}1$ cuboids.

Now assume the claim after slice $S_{h-1}$, so the current list is the single accumulated cuboid $A_{h-1}$. Consider the processing of $S_h$. Every cuboid created from points of $S_h$ has the singleton ruler $\{b_0+s_0 h\}$ in dimension~0. Therefore it can merge with $A_{h-1}$ only if it already agrees with $A_{h-1}$ in dimensions $1,\ldots,d-1$, which means that its rulers in those dimensions are exactly $R_1,\ldots,R_{d-1}$. That happens only when the whole slice has collapsed to the single cuboid $S_h$. Until then, $A_{h-1}$ cannot interfere with the slice-local merges. The behavior inside $S_h$ is therefore exactly the greedy-plus-merge behavior for the $(d{-}1)$-dimensional cuboid $R_1 \times \cdots \times R_{d-1}$. By the induction hypothesis on~$d$, the cuboids inside $S_h$ collapse to the single cuboid $S_h$ while using at most $d{-}1$ cuboids. At that point $A_{h-1}$ and $S_h$ agree in dimensions $1,\ldots,d-1$ and are adjacent in dimension~0, so the backward merge combines them into $A_h$.

This proves the claim for every slice. After the last slice, $A_{|R_0|-1}=C$, so the algorithm reconstructs the input as one cuboid. During the processing of any slice there is at most one accumulated cuboid from earlier slices and at most $d{-}1$ cuboids inside the current slice, so the total number of cuboids present at any time is at most $d$.
\end{proof}

\subsection{Greedy-Plus-Combine Optimality}
\label{proof:gpc-optimality}

We prove that the greedy-plus-combine always produces a decomposition
with the minimum number of cuboids.

\begin{figure}[ht]
\centering
\iptCornerSteal
\caption{Corner-steal and repair in $d = 2$.  The optimal decomposition has $C^*_0 = [1,4){\times}[1,2)$ (column) and $C^*_1 = [4,6){\times}[1,3)$ (block).  (a)~The greedy steals $p_1 = \operatorname{glb}(C^*_1) = (4,1)$ into the column.  (b)~After processing the remaining points, the algorithm holds the extended column~$C_0$, the partial fibre~$C_1 = \{(4,2)\}$, and the complete second slice~$C_2 = [5,6){\times}[1,3)$.  The bounding box $K = \operatorname{bbox}(C_1 \cup C_2)$ has a single deficit point $(4,1) \in C_0$.  (c)~The combine replaces $(C_0, C_1, C_2)$ by $C^*_0 = C_0 \setminus K$ and $C^*_1 = K$, restoring the optimum.}
\Description{A three-panel schematic of the corner-steal phenomenon in two dimensions. The first panel shows the optimal two-cuboid decomposition. The second shows the greedy algorithm after one point has been absorbed into the wrong cuboid and identifies the deficit in the bounding box of the last two cuboids. The third panel shows the repaired optimal decomposition after the combine step.}
\label{figure:corner-steal}
\end{figure}

\begin{proof}[Proof of greedy-plus-combine optimality]
By strong induction on~$N$, with a secondary induction on~$d$.
Let $(C^*_0, \ldots, C^*_{\kappa_{\text{opt}}-1})$ be a minimum-cuboid decomposition of~$D$, chosen so that $C^*_0$ is as large as possible. We show that the greedy-plus-combine reconstructs $C^*_0$ and never modifies it afterward.

The argument has three stages. First, we show that the algorithm reconstructs the first optimal cuboid~$C^*_0$. Second, if the greedy scan extends past~$C^*_0$, we show that the overshoot can only be a corner-steal in~$\dim_0$ and that the combine step repairs it immediately. Third, once $C^*_0$ has appeared, we show that it remains stable: in the no-steal branch because canonical maximality rules out merges and lexicographic separation rules out combines, and in the post-repair branch because a strict gap in~$\dim_0$ isolates $C^*_0$ from all later cuboids.

\paragraph{Recovery of~$C^*_0$.}
Consider the prefix of the input consisting exactly of the points of~$C^*_0$. Since $C^*_0$ is one cuboid, the one-cuboid recovery result from Appendix~\ref{proof:one-cuboid-recovery} shows that greedy-plus-merge would finish this prefix as one cuboid. It remains to check that the extra combine step cannot change that outcome.

The one-cuboid recovery proof processes $C^*_0$ slice by slice in dimension~0. While one slice is being processed, the cuboid list consists of one accumulated cuboid~$A$ from earlier slices together with the cuboids currently created inside the active slice. A combine involving $A$ cannot fire, because the bounding box of the last two cuboids lies entirely in the current slice, whereas~$A$ lies in earlier slices, so the deficit is disjoint from~$A$. A combine using only cuboids from the current slice also stays inside that slice. After deleting the fixed $\dim_0$ coordinate, it becomes exactly a greedy-plus-combine run on the $(d{-}1)$-dimensional one-cuboid input defined by that slice. By the secondary induction on~$d$, that lower-dimensional run reconstructs the slice as one cuboid. Hence each slice is finished correctly, and adjacent finished slices merge exactly as in one-cuboid recovery. When the last point of~$C^*_0$ has been processed, the current first cuboid is $C^*_0$.

If the greedy scan stops there, then $C_0 = C^*_0$ and we can pass directly to the stability argument below. Otherwise the greedy absorbs at least one point of the next optimal cuboid~$C^*_1$.

\paragraph{If the next point is absorbed, the direction must be $\dim_0$.}
Let $p_1 = \operatorname{glb}(C^*_1)$. For $C^*_0 \cup \{p_1\}$ to be a cuboid, the new point can extend exactly one ruler of~$C^*_0$. Therefore every other ruler of~$C^*_0$ has length~$1$, so $C^*_0$ is a column in some dimension~$m$. Moreover $C^*_1$ cannot be a column in that same dimension with the same fixed coordinates in all other dimensions, because then $C^*_0 \cup C^*_1$ would itself be a cuboid and the chosen decomposition would not be optimal.

Suppose $m > 0$. Then $C^*_0$ and the absorbed point~$p_1$ lie in one fixed $\dim_0$ slice. Before the algorithm leaves that slice, every point that can interact with this absorption also lies in that same slice. After deleting the fixed $\dim_0$ coordinate, both the input prefix and the corresponding part of the optimal decomposition become a $(d{-}1)$-dimensional instance with the same absorption pattern. By the secondary induction on~$d$, greedy-plus-combine already resolves that lower-dimensional instance optimally, so this absorption cannot persist when $m > 0$. Therefore the only remaining case is $m = 0$.

\paragraph{Repair of the corner-steal.}
Now $C^*_0$ varies only in~$\dim_0$ and has one fixed coordinate in every faster dimension. Let $m' \geq 1$ be the fastest dimension in which $C^*_1$ has extent greater than~$1$. Such a dimension exists, because otherwise $C^*_1$ would also be a $\dim_0$ column with the same fixed coordinates as~$C^*_0$, and the two cuboids would merge.

Partition $C^*_1$ into \emph{fibres}, where a fibre is the set of points that agree in every coordinate except~$\dim_{m'}$. These fibres arrive consecutively in lexicographic order. The first fibre contains~$p_1$. Since $p_1$ has already been absorbed into~$C_0$, the greedy scan collects the remaining $n'_{m'} - 1$ points of that fibre into a cuboid~$C_1$. The next fibre is complete, so it forms a cuboid~$C_2$ of length $n'_{m'}$ in~$\dim_{m'}$. The two cuboids cannot merge, because their rulers in~$\dim_{m'}$ have different lengths.

Let $K = \operatorname{bbox}(C_1 \cup C_2)$. This cuboid contains exactly the two fibres together with the single missing point~$p_1$. Hence the deficit of~$K$ is precisely $\{p_1\}$, which lies in~$C_0$. Removing~$K$ from the extended column restores the original column, so $C_0 \setminus K = C^*_0$ is again a cuboid. Therefore the combine step replaces $(C_0, C_1, C_2)$ by $(C^*_0, K)$ and repairs the earlier absorption of~$p_1$ into the current cuboid as soon as the second fibre has been processed.

From that point on, $K$ is the correct beginning of~$C^*_1$. The remaining fibres in the current $\dim_0$ slice are complete and arrive consecutively, so one-cuboid recovery merges them into~$K$ one after another. When that first $\dim_0$ slice of~$C^*_1$ is complete, it has become one cuboid. Every later $\dim_0$ slice of~$C^*_1$ is itself a complete $(d{-}1)$-dimensional cuboid. Again by one-cuboid recovery, each such slice collapses and then merges with the accumulated part of~$C^*_1$. Thus the algorithm reconstructs $C^*_1$ while keeping $C^*_0$ in place.

\paragraph{Stability of~$C^*_0$.}
Once the algorithm has a cuboid equal to~$C^*_0$, that cuboid is never modified again.  We treat the two possible configurations at the moment $C^*_0$ becomes the first cuboid separately, since they admit very different stability arguments: in the no-steal branch~$C^*_0$ may have arbitrary shape and the argument has to be carried out across all~$d$ dimensions, whereas after a corner-steal repair~$C^*_0$ is by construction a $\dim_0$-column whose $\dim_0$-range is strictly below every subsequent input point.

\emph{No-steal branch.}
If no steal occurs, $C_0 = C^*_0$ once the prefix corresponding to~$C^*_0$ has been processed, and the next input point is $p_1 = \operatorname{glb}(C^*_1)$. We rule out merge and combine separately: merge is excluded by the canonical-decomposition choice of~$C^*_0$, while combine is excluded by lexicographic separation.

\emph{No merge.}
Recall that the optimal decomposition $(C^*_0, \ldots, C^*_{\kappa_{\mathrm{opt}}-1})$ was chosen, among all minimum-$\kappa$ decompositions of~$D$, to maximise $|C^*_0|$ (canonical decomposition).  Suppose a later cuboid~$C_l$ ($l \geq 1$) merges with~$C^*_0$ during the run, so that $C^*_0 \cup C_l$ is a cuboid.  By the strong induction hypothesis on~$N$ applied to $D \setminus (C^*_0 \cup C_l)$, greedy-plus-combine produces a minimum-cuboid decomposition $(C''_0, \ldots)$ of that residual.  Then $(C^*_0 \cup C_l, C''_0, \ldots)$ is a minimum-cuboid decomposition of~$D$ whose first cuboid strictly exceeds~$C^*_0$ in size, contradicting the canonical-decomposition hypothesis.

\emph{No combine.}
Since $C^*_0$ is the first cuboid in the output list and the combine acts on the last three cuboids, a combine targeting~$C^*_0$ is necessarily the first combine, with triple $(C^*_0, C', C'')$ for two later cuboids~$C', C''$.  For it to fire, the bounding box $K = \operatorname{bbox}(C' \cup C'')$ must satisfy $K \cap C^*_0 \neq \emptyset$ (the deficit lies in~$C^*_0$).  We show $K \cap C^*_0 = \emptyset$ by induction on the dimension index.

Strict lexicographic ascent of the output cuboid list gives $\operatorname{lub}(C^*_0) <_{\mathrm{lex}} \operatorname{glb}(C') <_{\mathrm{lex}} \operatorname{glb}(C'')$.  Let $j \in \{0, \ldots, d{-}1\}$ be the smallest index with $\operatorname{lub}_j(C^*_0) < \operatorname{glb}_j(C')$, and let $j' \in \{0, \ldots, d{-}1\}$ be the smallest index with $\operatorname{glb}_{j'}(C') < \operatorname{glb}_{j'}(C'')$.  Suppose for contradiction that $q \in K \cap C^*_0$, so $q_i \leq \operatorname{lub}_i(C^*_0)$ and $q_i \geq \operatorname{glb}_i(K) = \min(\operatorname{glb}_i(C'), \operatorname{glb}_i(C''))$ for every~$i$.

\emph{Coordinates $i < \min(j, j')$.}
By definition of~$j$ and~$j'$, $\operatorname{glb}_i(C') = \operatorname{lub}_i(C^*_0)$ and $\operatorname{glb}_i(C'') = \operatorname{glb}_i(C')$, so $\operatorname{glb}_i(K) = \operatorname{lub}_i(C^*_0)$.  Combined with $q_i \leq \operatorname{lub}_i(C^*_0)$ this forces $q_i = \operatorname{lub}_i(C^*_0)$.

\emph{Case $j \leq j'$.}
At dimension~$j$ we have $\operatorname{glb}_j(C') > \operatorname{lub}_j(C^*_0)$ by definition of~$j$, and $\operatorname{glb}_j(C'') \geq \operatorname{glb}_j(C')$ because $j \leq j'$.  Hence $\operatorname{glb}_j(K) > \operatorname{lub}_j(C^*_0) \geq q_j$, contradicting $q_j \geq \operatorname{glb}_j(K)$.

\emph{Case $j' < j$.}
At dimension~$j'$, $\operatorname{glb}_{j'}(C') = \operatorname{lub}_{j'}(C^*_0)$ (since $j' < j$) and $\operatorname{glb}_{j'}(C'') > \operatorname{glb}_{j'}(C')$, so $\operatorname{glb}_{j'}(K) = \operatorname{lub}_{j'}(C^*_0)$ and the previous-paragraph argument forces $q_{j'} = \operatorname{lub}_{j'}(C^*_0) = \operatorname{glb}_{j'}(C')$.  In particular $q_{j'} < \operatorname{glb}_{j'}(C'')$, so every point of~$C''$ has $\dim_{j'}$-coordinate strictly greater than~$q_{j'}$ and $C''$ contributes nothing to the bounding box of~$K$ on the slice $\{x : x_{j'} = q_{j'}\}$.  Restricted to that slice, $K$ coincides with $\operatorname{bbox}(C')$.  Delete the first $j'{+}1$ coordinates and project to the remaining coordinates $\{j'{+}1, \ldots, d{-}1\}$.  The projected point still lies in both the projection of~$C^*_0$ and the projection of~$K$, and the first lexicographic disagreement between the projected cuboids now occurs at coordinate $j - j' - 1 \geq 0$.  Thus the same obstruction would reappear in $d - j' - 1$ dimensions for the projected pair $(C^*_0, C')$, contradicting the induction hypothesis for that lower-dimensional problem.

In either case $K \cap C^*_0 = \emptyset$, so no combine fires.

\emph{Post-repair branch.}
If the greedy scan absorbs the first point $p_1 = \operatorname{glb}(C^*_1)$ of the next optimal cuboid into the current first cuboid, then the repair above restores~$C^*_0$.  After the repair, $C^*_0$ varies only in~$\dim_0$ and $\operatorname{lub}_0(C^*_0) < p_{1,0}$.  Every later cuboid consists of points from $C^*_1 \cup C^*_2 \cup \cdots$ with $\dim_0$-coordinate $\geq p_{1,0}$.  Therefore any bounding box $K' = \operatorname{bbox}(C_{L-1} \cup C_L)$ formed from later cuboids satisfies $\operatorname{glb}_0(K') \geq p_{1,0} > \operatorname{lub}_0(C^*_0)$, making $K' \cap C^*_0 = \emptyset$ on the strength of $\dim_0$ alone, without invoking the multi-dimensional descent of the no-steal branch.  The same one-dimensional gap also rules out a merge: any merger with~$C^*_0$ would require an adjacent ruler in~$\dim_0$, but the smallest available $\dim_0$-coordinate among later points is~$p_{1,0}$, which strictly exceeds $\operatorname{lub}_0(C^*_0)$ rather than being adjacent to it.

\paragraph{Induction on~$N$.}
Since $C_0 = C^*_0$ and is stable, the algorithm processes $D' = D \setminus C^*_0$ independently.  By strong induction on~$N$ ($|D'| < |D|$), the greedy-plus-combine produces a minimum-cuboid decomposition of~$D'$.  The total cuboid count is $1 + (\kappa_{\text{opt}} - 1) = \kappa_{\text{opt}}$.
\end{proof}

\bibliographystyle{ACM-Reference-Format}
\bibliography{IPT}

\end{document}
